% =============================================================================
% Chapter 4 - Methodology
% =============================================================================
% Included from main.tex via % =============================================================================
% Chapter 4 - Methodology
% =============================================================================
% Included from main.tex via % =============================================================================
% Chapter 4 - Methodology
% =============================================================================
% Included from main.tex via % =============================================================================
% Chapter 4 - Methodology
% =============================================================================
% Included from main.tex via \include{contents/chapter4_shortened}. This file
% is the single source for Chapter 4; chapter4.tex and the chapter4_short_half
% files are superseded and share labels with it, so never include them.
%
% Rewritten 23 September 2026:
%   * Subsections consolidated; merged subsections keep every earlier label
%     (stacked \label commands), so references from other chapters resolve.
%   * Removed: "What is stored" and "Earlier evaluations" (labels
%     sec:what-is-stored, sec:excluded-evaluations, tab:five-evaluations),
%     all VRAM figures, the Computational Cost section (removed at the
%     author's request), and every code file, script and class name.
%   * The pipeline overview (formerly its own section) is now Section 4.4.1.
%   * New Section 4.8 combines the research gap and the limitations of each
%     component, grouped into motion modelling and temporal modelling.
%   * Figures fig4_1 ... fig4_4 are generated by tools/thesis_ch4_figures.py
%     from the values stated in this chapter.
%
% Requires in the preamble: amsmath, amssymb (\checkmark), graphicx, tikz with
%   \usetikzlibrary{arrows.meta, positioning, fit, backgrounds}
% =============================================================================

\chapter{Methodology}
\label{ch:methodology}

This chapter describes the study as it was carried out. Section~\ref{sec:research-approach} presents the experimental design, Section~\ref{sec:corpus-preparation} the dataset and its preparation, and Section~\ref{sec:ablation-matrix} the twelve configurations and the matched pairs used to compare them. Section~\ref{sec:architecture} describes the model architecture, Section~\ref{sec:training-procedure} the training procedure, and Section~\ref{sec:evaluation-protocol} the evaluation protocol. Section~\ref{sec:gaps-limitations} summarises the research gaps and limitations of each component, grouped into motion modelling and temporal modelling.


% =============================================================================
\section{Research Approach and Experimental Design}
\label{sec:research-approach}
% =============================================================================

\subsection{Factorial Ablation Design}
\label{sec:study-form}

This thesis does not propose and evaluate a single architecture. The study is designed as a factorial ablation: four components of a micro-expression recognition pipeline are each switched on or off, every valid combination is trained and evaluated under an identical protocol, and the contribution of each component is estimated from pairs of configurations that differ in that component alone. The object of study is therefore the set of configurations, and the question is not how well the pipeline performs but which of its stages are responsible for that performance, and by how much.

This design addresses the research gap of Section~\ref{sec:research-gap}. In the reviewed literature, components are varied one at a time within a single proposed system, for example the magnification factor~\cite{liX2018}, individual modules of a recurrent convolutional network~\cite{xia2020b} or the temporal aggregator of a transformer~\cite{zhang2022}. No reviewed study varies several components jointly under one protocol with all other factors held identical.

\subsection{Ablated Components}
\label{sec:four-variables}

The design space is defined by four components, each of which is either fully present or absent:

\begin{itemize}
	\item \textbf{Eulerian video magnification (EVM)}~\cite{wu2012}, applied during preprocessing before the motion representation is computed (Section~\ref{sec:motion-magnification}).
	\item \textbf{A convolutional spatial stem}, with one stream per motion channel (Section~\ref{sec:conv3d}).
	\item \textbf{SimAM}~\cite{yang2021}, a parameter-free attention module within the convolutional streams (Section~\ref{sec:simam}).
	\item \textbf{A temporal transformer encoder}~\cite{vaswani2017}, applied to the sequence of frame features (Section~\ref{sec:self-attention}).
\end{itemize}

The first three act on the motion representation and are grouped as \textit{motion modelling}; the transformer models the evolution of that representation over time and constitutes \textit{temporal modelling}. A disabled component is replaced by a minimal operation (Section~\ref{sec:fallbacks-config}) rather than by an alternative component of comparable capacity.

\subsection{Design Limitations}
\label{sec:design-limitation}
\label{sec:what-is-stored}

Every configuration is trained once, under a single fixed random seed (42). No configuration is retrained with other seeds, so no variance across seeds is estimated. Every difference between two configurations, including the matched-pair effects of Section~\ref{sec:matched-pairs}, is therefore a difference between two point estimates without a confidence interval, and it could change under a different seed. In addition, per-clip predictions were not retained; only aggregate metrics and confusion matrices were stored. A paired significance test such as McNemar's test~\cite{mcnemar1947}, which requires the prediction for each clip, therefore cannot be computed for any pair of configurations.


% =============================================================================
\section{Dataset and Preparation}
\label{sec:corpus-preparation}
% =============================================================================

\subsection{Data Selection}
\label{sec:corpus-used}
\label{sec:255-to-156}

All experiments use CASME~II~\cite{yan2014} (Section~\ref{sec:corpus}); no other dataset contributes to training or evaluation, so every result is a within-dataset result. The released coding contains 255 clips coded as micro-expressions with extractable frames (Section~\ref{sec:clip-selection}). The \texttt{Others} category (99 clips) is excluded before training. It is a residual category for clips that do not fit the other labels, and a class defined by exclusion rather than content cannot provide a meaningful training signal. Figure~\ref{fig:class-distribution} shows the composition of the 255 clips and the resulting working set of 156 clips.

\begin{figure}[htbp]
	\centering
	\includegraphics[width=0.85\textwidth]{fig4_1_class_distribution.png}
	\caption[Composition of the CASME II working set]{Composition of the 255 CASME~II micro-expression clips by emotion label, grouped into the three classes used in this thesis. The \texttt{Others} category is excluded, leaving a working set of 156 clips. Counts are taken from the released coding (Section~\ref{sec:label-taxonomy}).}
	\label{fig:class-distribution}
\end{figure}

\subsection{Class Grouping}
\label{sec:three-class-grouping}

Disgust, repression, sadness and fear are grouped into Negative, happiness forms Positive and surprise forms Surprise, following the valence-based convention of SMIC and MEGC 2019~\cite{liX2013,see2019} (Section~\ref{sec:label-taxonomy}). The resulting class distribution is 99 : 32 : 25, approximately 4 : 1.3 : 1. All results are reported for these three classes.

\subsection{Subjects and Folds}
\label{sec:subjects-loso}

The 255 clips come from 26 subjects. Subject 18 contributes only \texttt{Others} clips and is removed together with that category, so the working set contains 25 subjects, and leave-one-subject-out cross-validation produces \textbf{25 folds} rather than 26. Because subjects contribute unequal numbers of clips, the folds differ considerably in size and class composition (\autoref{fig:fold-composition}).

\begin{figure}[tbp]
	\centering
	\includegraphics[width=\textwidth]{fig3_2_fold_composition.png}
	\caption[Class composition of the 25 leave-one-subject-out folds]{Class composition of the 25 leave-one-subject-out folds. Top: clips held out in each fold, by class, with the number of classes present above each bar. Bottom: the number of folds containing one, two or three classes, and the per-fold macro-F1 ceiling this implies.}
	\label{fig:fold-composition}
\end{figure}

\subsection{Input Preparation}
\label{sec:network-preparation}

Each clip is converted into a motion tensor of shape $(3, 32, 224, 224)$, whose three channels hold horizontal optical flow, vertical optical flow and optical strain magnitude over 32 frame pairs (Sections~\ref{sec:recording-to-input} and~\ref{sec:video-to-motion}). Table~\ref{tab:preprocessing-settings} lists the settings used. None of them is varied by the ablation, so every configuration receives inputs of the same shape.

\begin{table}[htbp]
	\centering
	\small
	\begin{tabular}{@{}l l@{}}
		\hline
		\textbf{Setting} & \textbf{Value} \\
		\hline
		Frames sampled per clip        & 33, uniformly between onset and offset (32 flow pairs) \\
		Frame size                     & $224\times224$ pixels, greyscale \\
		Magnification factor $\alpha$  & 10 \\
		Magnification band             & 5--25 Hz at a nominal 200 fps \\
		Laplacian pyramid levels       & 4 \\
		Magnified output               & clipped to $[0,255]$ and converted to 8-bit \\
		Optical flow estimator         & Farneb\"ack~\cite{farneback2003} \\
		Channel normalisation          & min--max to $[0,1]$ per clip \\
		\hline
	\end{tabular}
	\caption[Preprocessing settings]{Preprocessing settings, fixed across all configurations. Section~\ref{sec:evm-departures} lists how the magnification settings depart from the reference formulation.}
	\label{tab:preprocessing-settings}
\end{table}


% =============================================================================
\section{Ablation Configurations}
\label{sec:ablation-matrix}
% =============================================================================

\subsection{Valid Configurations}
\label{sec:twelve-configs}
\label{sec:twelve-trained}
\label{sec:excluded-cells}
\label{sec:naming}

The four binary components define $2^4 = 16$ combinations (\autoref{fig:ablation-grid}). Four of them enable SimAM without the convolutional stem. Because SimAM re-weights the activations of a convolutional feature map, it has nothing to act on without the stem; these combinations are therefore not weaker models but ill-formed ones, and they are not trained. The remaining twelve configurations are listed in \autoref{tab:twelve-configs}. Configurations 1 to 8 were defined first and configurations 9, 12, 13 and 16 were added to complete the design, so the numbering omits the four excluded combinations (10, 11, 14 and 15). Configuration $n$ is denoted C$n$ in tables and figures.

\begin{figure}[htbp]
	\centering
	\includegraphics[width=0.8\textwidth]{fig4_2_ablation_grid.png}
	\caption[The factorial ablation design]{The $2^4$ factorial design. Each cell is one combination of the four components; twelve are trained and the four that enable SimAM without the convolutional stem are excluded.}
	\label{fig:ablation-grid}
\end{figure}

\begin{table}[htbp]
	\centering
	\small
	\begin{tabular}{@{}l c c c c@{}}
		\hline
		\textbf{Configuration} & \textbf{EVM} & \textbf{SimAM} & \textbf{CNN} & \textbf{TF} \\
		\hline
		C1   & --         & --         & --         & -- \\
		C2   & --         & --         & --         & \checkmark \\
		C3   & --         & --         & \checkmark & -- \\
		C4   & \checkmark & --         & --         & -- \\
		C5   & --         & \checkmark & \checkmark & -- \\
		C6   & --         & \checkmark & \checkmark & \checkmark \\
		C7   & \checkmark & --         & \checkmark & \checkmark \\
		C8   & \checkmark & \checkmark & \checkmark & \checkmark \\
		C9   & --         & --         & \checkmark & \checkmark \\
		C12  & \checkmark & --         & --         & \checkmark \\
		C13  & \checkmark & --         & \checkmark & -- \\
		C16  & \checkmark & \checkmark & \checkmark & -- \\
		\hline
	\end{tabular}
	\caption[The twelve valid ablation configurations]{The twelve valid ablation configurations. EVM: Eulerian video magnification; CNN: convolutional stem; TF: temporal transformer.}
	\label{tab:twelve-configs}
\end{table}

\subsection{Baseline and Control}
\label{sec:baseline-control}

Two configurations serve different purposes and should not be confused. C4 is the baseline of the study: it applies magnification and no learned component, and the contributions of the learned components are ultimately read against it. C1 is not a second baseline but the control for magnification: it is identical to C4 except that magnification is disabled. C8 enables all four components but is otherwise one of twelve configurations and is not privileged in the analysis.

\subsection{Matched Pairs}
\label{sec:matched-pairs}
\label{sec:matched-pairs-table}

A matched pair consists of two configurations that differ in exactly one component. The difference between their scores estimates the contribution of that component with all others held fixed, and the effect reported for a component is the mean over all of its pairs. Magnification and the transformer can be toggled in each of the six valid combinations of the other three components, giving six pairs each. A SimAM pair requires the stem to be present, and a stem pair requires SimAM to be absent, so these two components have four pairs each. Table~\ref{tab:matched-pairs} lists all twenty pairs. Because an effect averaged over four pairs is a noisier estimate than one averaged over six, effects with different pair counts should be compared with care.

\begin{table}[htbp]
	\centering
	\small
	\begin{tabular}{@{}l l r@{}}
		\hline
		\textbf{Component} & \textbf{Pairs (off $\rightarrow$ on)} & \textbf{Count} \\
		\hline
		Transformer & 1$\rightarrow$2, 4$\rightarrow$12, 3$\rightarrow$9, 13$\rightarrow$7, 5$\rightarrow$6, 16$\rightarrow$8 & 6 \\
		EVM         & 1$\rightarrow$4, 3$\rightarrow$13, 9$\rightarrow$7, 6$\rightarrow$8, 5$\rightarrow$16, 2$\rightarrow$12 & 6 \\
		SimAM       & 3$\rightarrow$5, 9$\rightarrow$6, 7$\rightarrow$8, 13$\rightarrow$16 & 4 \\
		CNN stem    & 1$\rightarrow$3, 2$\rightarrow$9, 4$\rightarrow$13, 12$\rightarrow$7 & 4 \\
		\hline
	\end{tabular}
	\caption[Matched pairs by component]{Matched pairs by component, identified by configuration number.}
	\label{tab:matched-pairs}
\end{table}


% =============================================================================
\section{Model Architecture}
\label{sec:architecture}
% =============================================================================

All twelve configurations are branches of a single model in which each component is executed or skipped according to switches fixed at construction; magnification is applied during preprocessing and is not part of the network. Because every configuration shares the same implementation, initialisation and layer order, two configurations that differ in one switch differ in nothing else.

\subsection{Pipeline Overview}
\label{sec:pipeline-overview}
\label{sec:pipeline-in-words}
\label{sec:pipeline-shapes-discussion}

The pipeline consists of two stages that are executed a different number of times (\autoref{fig:pipeline-overview}). \textbf{Stage~1 (preprocessing)} converts each clip into one motion tensor and runs once per clip. \textbf{Stage~2 (training and evaluation)} loads these tensors and fits a model; it runs once for every configuration and every held-out subject. Because Stage~1 is completed before any model is trained, all configurations are compared on identical inputs.

% -----------------------------------------------------------------------------
\begin{figure}[htbp]
	\centering
	\begin{tikzpicture}[
		font=\scriptsize,
		node distance=3.2mm,
		box/.style={
			draw, rounded corners=2pt, align=center,
			text width=34mm, inner sep=3pt,
			fill=black!4},
		switch/.style={box, dashed, fill=black!10},
		term/.style={box, fill=black!14, font=\scriptsize\bfseries},
		arr/.style={-{Stealth[length=1.6mm]}, semithick}
	]

	% ---------------- Stage 1: preprocessing ----------------
	\node[term] (a1) {CASME II dataset\\\normalfont 200\,fps, FACS-coded};
	\node[box, below=of a1] (a2)
		{Clip selection\\\normalfont 156 clips, 25 subjects};
	\node[box, below=of a2] (a3)
		{Frame preparation\\\normalfont greyscale, $224\times224$};
	\node[box, below=of a3] (a4)
		{Uniform sampling\\\normalfont 33 frames per clip};
	\node[switch, below=of a4] (a5)
		{Motion magnification\\\normalfont Eulerian, $\alpha=10$};
	\node[box, below=of a5] (a6)
		{Optical flow\\\normalfont Farneb\"ack, 32 pairs};
	\node[box, below=of a6] (a7)
		{Optical strain\\\normalfont spatial gradients};
	\node[box, below=of a7] (a8)
		{Stack and normalise\\\normalfont $[u,v,\varepsilon]$ to $[0,1]$};
	\node[term, below=of a8] (a9)
		{Saved motion tensor\\\normalfont $(3,32,224,224)$};

	\foreach \i/\j in {a1/a2, a2/a3, a3/a4, a4/a5, a5/a6, a6/a7, a7/a8, a8/a9}
		{\draw[arr] (\i) -- (\j);}

	% ---------------- Stage 2: training and evaluation ----------------
	\node[box, right=20mm of a2] (b1)
		{Load and standardise\\\normalfont flip augmentation};
	\node[switch, below=of b1] (b2)
		{Convolutional stem\\\normalfont three streams};
	\node[switch, below=of b2] (b3)
		{SimAM\\\normalfont parameter-free gate};
	\node[box, below=of b3] (b4)
		{Pool to a sequence\\\normalfont 32 steps, 96 features};
	\node[switch, below=of b4] (b5)
		{Temporal transformer\\\normalfont 4 layers, 8 heads};
	\node[box, below=of b5] (b6)
		{Classifier head\\\normalfont three logits};
	\node[box, below=of b6] (b7)
		{Leave-one-subject-out\\\normalfont 25 folds, counts pooled};
	\node[term, below=of b7] (b8)
		{Pooled macro F1\\\normalfont the primary metric};

	\foreach \i/\j in {b1/b2, b2/b3, b3/b4, b4/b5, b5/b6, b6/b7, b7/b8}
		{\draw[arr] (\i) -- (\j);}

	% ---------------- the hand-off ----------------
	\draw[arr] (a9.east) -- ++(8mm,0) |- (b1.west);

	% ---------------- stage brackets ----------------
	\begin{scope}[on background layer]
		\node[draw=black!45, dashed, rounded corners=3pt, inner sep=4.5mm,
		      fit=(a1)(a9)] (s1) {};
		\node[draw=black!45, dashed, rounded corners=3pt, inner sep=4.5mm,
		      fit=(b1)(b8)] (s2) {};
	\end{scope}
	\node[above=0.5mm of s1, font=\scriptsize\bfseries] {Stage 1: preprocessing};
	\node[above=0.5mm of s2, font=\scriptsize\bfseries] {Stage 2: training and evaluation};

	\end{tikzpicture}
	\caption[The data pipeline, end to end]{The pipeline from dataset to score. Stage~1 runs once per clip; Stage~2 runs once per configuration and fold. Solid boxes are always executed. \textbf{Dashed boxes are the four components switched on and off in the ablation}; each is replaced by a defined operation when disabled.}
	\label{fig:pipeline-overview}
\end{figure}
% -----------------------------------------------------------------------------

In Stage~1, each selected clip is loaded in greyscale, resized to $224\times224$ pixels and sampled at 33 evenly spaced frames between onset and offset. Magnification is applied at this point if enabled. Optical flow is then computed for the 32 adjacent frame pairs, optical strain is derived from its spatial gradients, and the three fields are stacked and normalised into one tensor. The network therefore never receives pixel intensities, only displacement and deformation fields. In Stage~2, each tensor is standardised per channel, optionally flipped for augmentation, and passed through the convolutional stem, the SimAM gate, a pooling step that removes the spatial dimensions, the temporal encoder and the classification head. Table~\ref{tab:pipeline-shapes} lists the representation of the data after each step.

\begin{table}[htbp]
	\centering
	\small
	\begin{tabular}{@{}lll@{}}
		\hline
		\textbf{After this step} & \textbf{The data is} & \textbf{Shape} \\
		\hline
		Clip selection        & a set of clips                & 156 clips, 25 subjects \\
		Frame preparation     & greyscale frames              & $N \times 224 \times 224$ \\
		Uniform sampling      & a fixed-length sequence       & $33 \times 224 \times 224$ \\
		Magnification         & the same, amplified           & $33 \times 224 \times 224$ \\
		Flow and strain       & three motion fields per pair  & $3 \times 32 \times 224 \times 224$ \\
		Pooling to a sequence & a sequence of feature vectors & $32 \times 96$ \\
		Temporal encoding     & one vector per clip           & $96$ \\
		Classifier head       & class scores                  & $3$ \\
		Fold aggregation      & pooled confusion counts       & $3 \times 3$ \\
		\hline
	\end{tabular}
	\caption[What the data is at each pipeline step]{The data representation after each pipeline step, where $N$ is the number of frames between onset and offset. The temporal axis falls from 33 frames to 32 pairs because flow is computed between \emph{adjacent} frames, and the spatial axes are removed at the pooling step rather than at the classifier.}
	\label{tab:pipeline-shapes}
\end{table}

\subsection{Convolutional Stem}
\label{sec:input-streams}
\label{sec:conv-stem}

Each of the three motion channels is processed by its own convolutional stream with unshared weights. Table~\ref{tab:conv-stem} lists the layers of one stream, which use batch normalisation~\cite{ioffe2015}, ReLU activations~\cite{nair2010} and channel-wise dropout~\cite{tompson2015}. No kernel spans the temporal axis (Section~\ref{sec:conv3d}), so the 32 time steps pass through unchanged while height and width are halved once. The outputs of the three streams are concatenated along the channel dimension, giving a feature map of $96\times32\times112\times112$ per clip.

\begin{table}[htbp]
	\centering
	\small
	\begin{tabular}{@{}l l l@{}}
		\hline
		\textbf{Layer} & \textbf{Configuration} & \textbf{Output (per stream)} \\
		\hline
		3D convolution          & 16 filters, kernel $(1,3,3)$, padding $(0,1,1)$, no bias & $16\times32\times224\times224$ \\
		Normalisation and activation & batch normalisation, ReLU, 3D dropout ($p=0.3$) & $16\times32\times224\times224$ \\
		3D convolution          & 32 filters, kernel $(1,3,3)$, padding $(0,1,1)$, no bias & $32\times32\times224\times224$ \\
		Normalisation and activation & batch normalisation, ReLU, 3D dropout ($p=0.3$) & $32\times32\times224\times224$ \\
		3D max pooling          & kernel and stride $(1,2,2)$ & $32\times32\times112\times112$ \\
		\hline
	\end{tabular}
	\caption[Layers of one convolutional stream]{Layers of one convolutional stream. Shapes are given as channels $\times$ time $\times$ height $\times$ width.}
	\label{tab:conv-stem}
\end{table}

\subsection{SimAM}
\label{sec:simam-placement}

When enabled, SimAM~\cite{yang2021} is applied separately to the output of each convolutional stream, before the streams are concatenated. It follows \eqref{eq:simam-stats} and \eqref{eq:simam-gate}, with the mean and variance computed over the time, height and width of each channel (divisor $THW-1$) and $\lambda = 10^{-4}$. It adds no learnable parameters. SimAM is enabled only together with the convolutional stem.

\subsection{Temporal Transformer}
\label{sec:transformer-config}

The concatenated feature map is reduced by adaptive average pooling over the two spatial dimensions to a sequence of 32 time steps, each a 96-dimensional feature vector. Fixed sinusoidal positional encodings~\cite{vaswani2017} are added to this sequence:
\begin{equation}
	\mathrm{PE}(p, 2i) = \sin\!\left(\frac{p}{10000^{2i/96}}\right), \qquad \mathrm{PE}(p, 2i+1) = \cos\!\left(\frac{p}{10000^{2i/96}}\right),
	\label{eq:positional-encoding}
\end{equation}
where $\mathrm{PE}(p, j)$ is the value added to feature $j$ of the token at time step $p \in \{0, \dots, 31\}$, $i$ indexes the pairs of feature dimensions, and 96 is the feature dimension. The encoding is fixed rather than learned, and dropout with probability 0.1 is applied after it is added. The sequence is then processed by a transformer encoder with four layers, eight attention heads, model dimension 96, feed-forward dimension 256, dropout 0.1, GELU activation~\cite{hendrycks2016} and pre-norm layer ordering~\cite{xiong2020}, followed by a final layer normalisation~\cite{ba2016}. The output sequence is averaged over time, giving one 96-dimensional vector per clip.

\subsection{Classification Head and Replacement Operations}
\label{sec:classifier-head}
\label{sec:fallbacks-config}

Every configuration ends in the same classification head: layer normalisation, dropout with probability 0.3, and a linear layer mapping the 96 features to three class logits. A disabled component is replaced by a minimal operation rather than removed, so that the configuration remains a fair control (\autoref{tab:ch2-component-controls}). Without the convolutional stem, each frame of the input tensor is average-pooled to a $4\times4$ grid per channel, and the resulting 48 values are projected to 96 features by a linear layer with 4,704 parameters. Without the transformer, the sequence is averaged over time; this operation has no parameters and no notion of temporal order.

\subsection{Parameter Counts}
\label{sec:parameter-counts}
\label{sec:weight-init}

The convolutional stem has 14,544 parameters across its three streams, the transformer encoder 348,736, the replacement projection 4,704 and the classification head 483; SimAM and magnification add none. Figure~\ref{fig:parameter-counts} shows the resulting totals for the four architectural combinations. In a configuration with all components, the transformer holds approximately 96\,\% of the parameters. All configurations use the same initialisation: Kaiming-normal initialisation for convolutional weights, Xavier-uniform initialisation for linear weights, unit scale and zero shift for normalisation layers, and zero biases.

\begin{figure}[htbp]
	\centering
	\includegraphics[width=0.85\textwidth]{fig4_3_parameter_counts.png}
	\caption[Parameter totals by architectural combination]{Trainable parameters of the four architectural combinations (logarithmic scale), with the contribution of each component.}
	\label{fig:parameter-counts}
\end{figure}


% =============================================================================
\section{Training Procedure}
\label{sec:training-procedure}
% =============================================================================

All twelve configurations are trained with the same procedure, summarised in \autoref{tab:training-settings}; the underlying methods are described in Section~\ref{sec:scarcity-skew}.

\begin{table}[htbp]
	\centering
	\small
	\begin{tabular}{@{}l p{0.62\linewidth}@{}}
		\hline
		\textbf{Setting} & \textbf{Value} \\
		\hline
		Loss            & focal loss~\cite{lin2017} with $\gamma = 2$~\cite{zhaoS2021} and label smoothing~\cite{szegedy2016} of 0.05; no class weights \\
		Sampling        & class-balanced, weight $1/n_c$ per example, with replacement \\
		Augmentation    & random horizontal flip ($p = 0.5$) with negation of the horizontal flow \\
		Optimiser       & AdamW~\cite{loshchilov2019}, learning rate $10^{-4}$, weight decay $10^{-4}$ \\
		Learning rate   & linear warmup over 5 epochs from $10^{-5}$, then cosine annealing~\cite{loshchilov2017} to $10^{-7}$ \\
		Epochs          & 50, without early stopping \\
		Batch size      & 8 \\
		Gradient clipping & global norm 1.0 \\
		Precision       & mixed precision with gradient scaling~\cite{micikevicius2018} \\
		Random seed     & 42, reset before every configuration and fold \\
		\hline
	\end{tabular}
	\caption[Training settings]{Training settings, identical for all configurations and folds.}
	\label{tab:training-settings}
\end{table}

\subsection{Loss and Class Weighting}
\label{sec:loss-config}
\label{sec:class-weighting-rule}
\label{sec:sampler-config}

The loss can apply per-class weights ($\alpha_t$ in \eqref{eq:focal-loss}), but they are enabled only when the class-balanced sampler is inactive. Because the sampler is active in every reported run, the loss carries \textbf{no class weights} (Section~\ref{sec:combining-both}), and the frequency correction comes entirely from the sampler, which gives each training example a weight inversely proportional to the size of its class in the training split of the current fold \eqref{eq:sampler-weight}. The focal term $(1-p_t)^\gamma$ and label smoothing remain active; neither depends on class frequency.

The rule was introduced after two failures in opposite directions, both observed during the development of this study. In an earlier version, the loss carried inverse-frequency class weights and no sampler was used; under this single correction, the proposed configuration stopped predicting the \textbf{Positive} class altogether. A later version added the sampler while keeping the loss weights, so both corrections operated together; the failure then reversed, and configurations stopped predicting the \textbf{Negative} majority class, one of them predicting a single class for every clip. With the rule in place, no configuration in the reported study abandons any class. The only reviewed study that combines two corrections reports that the balanced loss adds only a slight improvement, not that the combination reverses its effect~\cite{xia2020a}. The two failures, however, occurred in implementation versions that differ from the reported run in more than the correction rule, so \textbf{the causal effect of the rule is not isolated}.

\subsection{Optimisation}
\label{sec:optimiser-config}
\label{sec:batch-size}
\label{sec:numerical-safeguards}

Figure~\ref{fig:lr-schedule} shows the learning-rate schedule. The learning rate rises linearly from one tenth of its target value over the first five epochs and then follows a cosine decay to $10^{-7}$; the schedule is updated once per epoch. Every configuration and fold is trained for the full 50 epochs. Gradients are clipped to a global norm of 1.0 after being unscaled from mixed precision, and an update is skipped if any gradient contains a non-finite value.

\begin{figure}[htbp]
	\centering
	\includegraphics[width=0.8\textwidth]{fig4_4_lr_schedule.png}
	\caption[Learning-rate schedule]{Learning rate per epoch: linear warmup over five epochs followed by cosine annealing.}
	\label{fig:lr-schedule}
\end{figure}

\subsection{Reproducibility}
\label{sec:seeding}
\label{sec:no-hyperparameter-search}

The random seed is reset before every configuration and every fold, so the sampling order, weight initialisation and all other stochastic steps start from the same state each time a fold is trained. No hyperparameter search was performed: all training settings are fixed defaults, identical across the twelve configurations and the 25 folds. Every comparison is therefore between architectures trained with one common, untuned recipe, not between architectures each tuned to its own advantage.


% =============================================================================
\section{Evaluation Protocol}
\label{sec:evaluation-protocol}
% =============================================================================

All configurations are evaluated with the same fold structure, checkpoint selection rule and aggregation across folds. The underlying metrics are defined in Section~\ref{sec:evaluation}.

\subsection{Leave-One-Subject-Out Folds}
\label{sec:loso-folds}

Evaluation uses leave-one-subject-out cross-validation~\cite{yan2014,see2019} over the 25 subjects of the working set, giving \textbf{25 folds}. Each fold holds out all clips of one subject and trains on the clips of the remaining 24 subjects. All 25 folds are run for every configuration.

\subsection{Checkpoint Selection}
\label{sec:checkpoint-selection}

Within each fold, the macro F1 on the held-out subject is computed after every epoch, and the weights of the best epoch are restored and scored on the same held-out subject. \textbf{There is no inner validation split}: the held-out clips are used both to select the epoch and to score it. A proper nested procedure would select the epoch on a validation subset of the training subjects and reserve the held-out subject for the final evaluation. The reported scores are therefore optimistically biased by an unknown amount~\cite{varma2006}. The bias is likely to be largest for the smallest folds, where the choice among 50 epochs is made on only a few clips.

\subsection{Aggregation and Primary Metric}
\label{sec:aggregation}
\label{sec:primary-metric}
\label{sec:mean-of-folds-rejected}

The confusion matrices of the 25 folds are summed into one pooled $3\times3$ matrix, from which the per-class precision, recall and F1 are computed. The primary metric is \textbf{pooled macro F1}, the mean of the three pooled per-class F1 scores, which is identical in construction to the Unweighted F1 of MEGC 2019~\cite{see2019}. Pooled accuracy is reported alongside it.

The alternative, averaging the macro F1 of each fold, is rejected because of the fold composition. Of the 25 folds, 10 contain clips of only one class, 8 contain two and 7 contain all three. A class absent from a fold contributes an F1 of zero to that fold, so a single-class fold can score at most $1/3$ and a two-class fold at most $2/3$. The mean-of-folds macro F1 is therefore bounded by
\begin{equation}
	\overline{F1}_{\max} = \frac{10 \cdot \tfrac13 + 8 \cdot \tfrac23 + 7 \cdot 1}{25} = \frac{3.333\ldots + 5.333\ldots + 7}{25} \approx 0.627,
	\label{eq:fold-ceiling}
\end{equation}
where $\overline{F1}_{\max}$ is the highest mean-of-folds macro F1 that any configuration can reach, the numerators weight each fold type by its maximum score, and 25 is the number of folds. No configuration can exceed this ceiling, however accurate it is. Pooling avoids the ceiling, because true positives, false positives and false negatives are accumulated over all 156 clips before any F1 is computed. Fold composition also inflates mean-of-folds accuracy: for the configuration with all four components, it exceeds pooled accuracy by 6.3 percentage points.


% =============================================================================
\section{Research Gaps and Limitations}
\label{sec:gaps-limitations}
% =============================================================================

This section combines, for each component, the research gap left by the reviewed literature (Chapter~\ref{ch:literature}) and the limitations of its implementation in this thesis. The components are grouped into motion modelling, which covers the representation and spatial processing of facial motion, and temporal modelling, which covers the sampling of the sequence and the modelling of its evolution over time. Limitations of the training and evaluation procedure as a whole are listed last.

\subsection{Motion Modelling}
\label{sec:gaps-motion}

\subsubsection{Eulerian Video Magnification}

\textbf{Research gap.} Magnification improves recognition in the reviewed studies, but every evaluation places it in front of an appearance-based representation such as LBP, HOG or pre-trained face features~\cite{liX2018,bai2021,liY2018,liY2021}. No study measures its contribution when the classifier receives an explicit motion field. The magnification factor has an interior optimum that can be found only by search~\cite{liX2018}, and magnification also amplifies unwanted motion such as head movement~\cite{liX2018}.

\textbf{Limitation.} The factor is fixed at $\alpha = 10$, inside the reported optimum, but is not varied. The pass-band, the filter and the order of magnification and sampling depart from the reference formulation (\autoref{tab:ch2-evm-departures}); in particular, the realised pass-band depends on the clip length and is approximately 2.5--12.5~Hz at the median length. No face alignment precedes magnification, so head motion within the band is amplified together with facial movement.

\subsubsection{Optical Flow}

\textbf{Research gap.} Optical flow is the preferred input of the strongest benchmark systems~\cite{see2019}, but the choice of estimator is not settled: Farneb\"ack and TV-L1 are both used~\cite{zhaoS2021,liong2019a}, and no controlled comparison of their recognition performance exists. Whether a single onset-to-apex flow field~\cite{liong2019a,liong2019b} or a sequence of flow fields~\cite{zhaoS2021} is preferable has likewise not been measured.

\textbf{Limitation.} Only the Farneb\"ack estimator~\cite{farneback2003} is used, and flow is computed between consecutive sampled frames rather than relative to the onset. No flow-based fine alignment is performed, although a small misalignment can considerably degrade micro-expression recognition~\cite{xu2017}.

\subsubsection{Optical Strain}

\textbf{Research gap.} The evidence for strain comes from hand-crafted descriptors evaluated under a protocol that is not subject-independent~\cite{liong2016}, and the strain magnitude of STSTNet cannot be determined unambiguously from its published equation~\cite{liong2019b}.

\textbf{Limitation.} The implemented magnitude counts shear once rather than twice, weighting it differently from the Frobenius norm used by Liong et al.~\cite{liong2014a,liong2016}, and no dedicated strain filter is applied~\cite{liong2014a}. Strain is present in all twelve configurations, so its own contribution is not measured.

\subsubsection{Convolutional Spatial Stem}

\textbf{Research gap.} The reviewed studies add modules to a backbone that remains in place~\cite{xia2020b}; none removes the learned spatial stem and passes the motion representation directly to the classifier, so its contribution is unknown. Deeper models and higher input resolutions reduce performance on composite micro-expression data~\cite{xia2020b}.

\textbf{Limitation.} The stem contains only spatial kernels, so the ablation tests a learned spatial stem rather than spatio-temporal convolution~\cite{ji2013,tran2015}. Its input resolution of $224\times224$ exceeds the $100\times100$ threshold reported by Xia et al.~\cite{xia2020b}, so any effect measured for the stem may reflect resolution as much as architecture.

\subsubsection{Parameter-Free Attention}

\textbf{Research gap.} SimAM has been evaluated on natural images with learned appearance features~\cite{yang2021}; its suitability for flow-derived features, and whether its statistics should be computed per frame or per clip, are untested. Its regularisation constant was selected separately for each image dataset~\cite{yang2021}.

\textbf{Limitation.} SimAM statistics are computed over the whole clip rather than per frame, and $\lambda = 10^{-4}$ is not re-tuned. Because SimAM requires the convolutional stem, its effect is measured over only four matched pairs and is conditional on the presence of the stem.

\subsection{Temporal Modelling}
\label{sec:gaps-temporal}

\subsubsection{Temporal Normalisation}

\textbf{Research gap.} Sequence length is usually fixed at ten frames by convention~\cite{liX2013,lu2015}, and the studies that vary it report different optima~\cite{liX2018,ben2021,xu2017}. No reviewed study compares interpolation methods for a motion-based representation, in which synthesised frames would produce synthesised displacements; Xu et al. prefer linear interpolation for motion features but do not compare it with TIM~\cite{xu2017}.

\textbf{Limitation.} Temporal normalisation is implemented as uniform sampling of existing frames rather than the Temporal Interpolation Model~\cite{pfister2011}, so short clips cannot be lengthened and three clips repeat frames. The length of 33 frames is not varied. Sampling reduces the effective frame rate to approximately 100~fps at the median clip length and 52~fps for the longest clip (Section~\ref{sec:three-channel-tensor}).

\subsubsection{Temporal Transformer}

\textbf{Research gap.} The strongest transformer-based results combine a pre-trained spatial transformer with an LSTM aggregator~\cite{zhang2022}, so the benefit of a temporal transformer trained from scratch on one small dataset has not been isolated. Transformers lack the inductive biases of convolution and can perform poorly when trained on insufficient data~\cite{dosovitskiy2021}.

\textbf{Limitation.} The implemented encoder differs from SLSTT in six respects (\autoref{tab:slstt-divergences}). The ablation compares the complete encoder, including positional encoding, feed-forward layers and normalisation, with mean pooling; it does not isolate self-attention or compare it with another learned temporal model.

\subsection{Training and Evaluation}
\label{sec:gaps-training}

\textbf{Research gap.} How imbalance corrections interact is examined by only one reviewed study~\cite{xia2020a}, the focusing parameter of focal loss has not been validated for this task~\cite{zhaoS2021}, and label smoothing~\cite{szegedy2016} has not been examined for micro-expression recognition.

\textbf{Limitation.} The imbalance treatment, $\gamma$ and label smoothing are identical in all configurations, so their effects are not measured. All results come from one dataset, one random seed and a checkpoint selection without an inner validation split (Sections~\ref{sec:design-limitation} and~\ref{sec:checkpoint-selection}), and no result extends to other datasets, populations or recording conditions (Section~\ref{sec:sibling-corpora}).


% =============================================================================
\section{Summary}
\label{sec:ch4-conclusion}
% =============================================================================

This chapter described the factorial ablation of four pipeline components on CASME~II. It presented the 156-clip working set and its 25 leave-one-subject-out folds, the twelve valid configurations and the twenty matched pairs used to compare them, and the shared model architecture, training procedure and evaluation protocol, with pooled macro F1 as the primary metric. It closed with the research gaps and limitations of each motion-modelling and temporal-modelling component.

Chapter~\ref{ch:results} reports the results of the ablation: the performance of the twelve configurations and the contribution of each component, measured over its matched pairs.
. This file
% is the single source for Chapter 4; chapter4.tex and the chapter4_short_half
% files are superseded and share labels with it, so never include them.
%
% Rewritten 23 September 2026:
%   * Subsections consolidated; merged subsections keep every earlier label
%     (stacked \label commands), so references from other chapters resolve.
%   * Removed: "What is stored" and "Earlier evaluations" (labels
%     sec:what-is-stored, sec:excluded-evaluations, tab:five-evaluations),
%     all VRAM figures, the Computational Cost section (removed at the
%     author's request), and every code file, script and class name.
%   * The pipeline overview (formerly its own section) is now Section 4.4.1.
%   * New Section 4.8 combines the research gap and the limitations of each
%     component, grouped into motion modelling and temporal modelling.
%   * Figures fig4_1 ... fig4_4 are generated by tools/thesis_ch4_figures.py
%     from the values stated in this chapter.
%
% Requires in the preamble: amsmath, amssymb (\checkmark), graphicx, tikz with
%   \usetikzlibrary{arrows.meta, positioning, fit, backgrounds}
% =============================================================================

\chapter{Methodology}
\label{ch:methodology}

This chapter describes the study as it was carried out. Section~\ref{sec:research-approach} presents the experimental design, Section~\ref{sec:corpus-preparation} the dataset and its preparation, and Section~\ref{sec:ablation-matrix} the twelve configurations and the matched pairs used to compare them. Section~\ref{sec:architecture} describes the model architecture, Section~\ref{sec:training-procedure} the training procedure, and Section~\ref{sec:evaluation-protocol} the evaluation protocol. Section~\ref{sec:gaps-limitations} summarises the research gaps and limitations of each component, grouped into motion modelling and temporal modelling.


% =============================================================================
\section{Research Approach and Experimental Design}
\label{sec:research-approach}
% =============================================================================

\subsection{Factorial Ablation Design}
\label{sec:study-form}

This thesis does not propose and evaluate a single architecture. The study is designed as a factorial ablation: four components of a micro-expression recognition pipeline are each switched on or off, every valid combination is trained and evaluated under an identical protocol, and the contribution of each component is estimated from pairs of configurations that differ in that component alone. The object of study is therefore the set of configurations, and the question is not how well the pipeline performs but which of its stages are responsible for that performance, and by how much.

This design addresses the research gap of Section~\ref{sec:research-gap}. In the reviewed literature, components are varied one at a time within a single proposed system, for example the magnification factor~\cite{liX2018}, individual modules of a recurrent convolutional network~\cite{xia2020b} or the temporal aggregator of a transformer~\cite{zhang2022}. No reviewed study varies several components jointly under one protocol with all other factors held identical.

\subsection{Ablated Components}
\label{sec:four-variables}

The design space is defined by four components, each of which is either fully present or absent:

\begin{itemize}
	\item \textbf{Eulerian video magnification (EVM)}~\cite{wu2012}, applied during preprocessing before the motion representation is computed (Section~\ref{sec:motion-magnification}).
	\item \textbf{A convolutional spatial stem}, with one stream per motion channel (Section~\ref{sec:conv3d}).
	\item \textbf{SimAM}~\cite{yang2021}, a parameter-free attention module within the convolutional streams (Section~\ref{sec:simam}).
	\item \textbf{A temporal transformer encoder}~\cite{vaswani2017}, applied to the sequence of frame features (Section~\ref{sec:self-attention}).
\end{itemize}

The first three act on the motion representation and are grouped as \textit{motion modelling}; the transformer models the evolution of that representation over time and constitutes \textit{temporal modelling}. A disabled component is replaced by a minimal operation (Section~\ref{sec:fallbacks-config}) rather than by an alternative component of comparable capacity.

\subsection{Design Limitations}
\label{sec:design-limitation}
\label{sec:what-is-stored}

Every configuration is trained once, under a single fixed random seed (42). No configuration is retrained with other seeds, so no variance across seeds is estimated. Every difference between two configurations, including the matched-pair effects of Section~\ref{sec:matched-pairs}, is therefore a difference between two point estimates without a confidence interval, and it could change under a different seed. In addition, per-clip predictions were not retained; only aggregate metrics and confusion matrices were stored. A paired significance test such as McNemar's test~\cite{mcnemar1947}, which requires the prediction for each clip, therefore cannot be computed for any pair of configurations.


% =============================================================================
\section{Dataset and Preparation}
\label{sec:corpus-preparation}
% =============================================================================

\subsection{Data Selection}
\label{sec:corpus-used}
\label{sec:255-to-156}

All experiments use CASME~II~\cite{yan2014} (Section~\ref{sec:corpus}); no other dataset contributes to training or evaluation, so every result is a within-dataset result. The released coding contains 255 clips coded as micro-expressions with extractable frames (Section~\ref{sec:clip-selection}). The \texttt{Others} category (99 clips) is excluded before training. It is a residual category for clips that do not fit the other labels, and a class defined by exclusion rather than content cannot provide a meaningful training signal. Figure~\ref{fig:class-distribution} shows the composition of the 255 clips and the resulting working set of 156 clips.

\begin{figure}[htbp]
	\centering
	\includegraphics[width=0.85\textwidth]{fig4_1_class_distribution.png}
	\caption[Composition of the CASME II working set]{Composition of the 255 CASME~II micro-expression clips by emotion label, grouped into the three classes used in this thesis. The \texttt{Others} category is excluded, leaving a working set of 156 clips. Counts are taken from the released coding (Section~\ref{sec:label-taxonomy}).}
	\label{fig:class-distribution}
\end{figure}

\subsection{Class Grouping}
\label{sec:three-class-grouping}

Disgust, repression, sadness and fear are grouped into Negative, happiness forms Positive and surprise forms Surprise, following the valence-based convention of SMIC and MEGC 2019~\cite{liX2013,see2019} (Section~\ref{sec:label-taxonomy}). The resulting class distribution is 99 : 32 : 25, approximately 4 : 1.3 : 1. All results are reported for these three classes.

\subsection{Subjects and Folds}
\label{sec:subjects-loso}

The 255 clips come from 26 subjects. Subject 18 contributes only \texttt{Others} clips and is removed together with that category, so the working set contains 25 subjects, and leave-one-subject-out cross-validation produces \textbf{25 folds} rather than 26. Because subjects contribute unequal numbers of clips, the folds differ considerably in size and class composition (\autoref{fig:fold-composition}).

\begin{figure}[tbp]
	\centering
	\includegraphics[width=\textwidth]{fig3_2_fold_composition.png}
	\caption[Class composition of the 25 leave-one-subject-out folds]{Class composition of the 25 leave-one-subject-out folds. Top: clips held out in each fold, by class, with the number of classes present above each bar. Bottom: the number of folds containing one, two or three classes, and the per-fold macro-F1 ceiling this implies.}
	\label{fig:fold-composition}
\end{figure}

\subsection{Input Preparation}
\label{sec:network-preparation}

Each clip is converted into a motion tensor of shape $(3, 32, 224, 224)$, whose three channels hold horizontal optical flow, vertical optical flow and optical strain magnitude over 32 frame pairs (Sections~\ref{sec:recording-to-input} and~\ref{sec:video-to-motion}). Table~\ref{tab:preprocessing-settings} lists the settings used. None of them is varied by the ablation, so every configuration receives inputs of the same shape.

\begin{table}[htbp]
	\centering
	\small
	\begin{tabular}{@{}l l@{}}
		\hline
		\textbf{Setting} & \textbf{Value} \\
		\hline
		Frames sampled per clip        & 33, uniformly between onset and offset (32 flow pairs) \\
		Frame size                     & $224\times224$ pixels, greyscale \\
		Magnification factor $\alpha$  & 10 \\
		Magnification band             & 5--25 Hz at a nominal 200 fps \\
		Laplacian pyramid levels       & 4 \\
		Magnified output               & clipped to $[0,255]$ and converted to 8-bit \\
		Optical flow estimator         & Farneb\"ack~\cite{farneback2003} \\
		Channel normalisation          & min--max to $[0,1]$ per clip \\
		\hline
	\end{tabular}
	\caption[Preprocessing settings]{Preprocessing settings, fixed across all configurations. Section~\ref{sec:evm-departures} lists how the magnification settings depart from the reference formulation.}
	\label{tab:preprocessing-settings}
\end{table}


% =============================================================================
\section{Ablation Configurations}
\label{sec:ablation-matrix}
% =============================================================================

\subsection{Valid Configurations}
\label{sec:twelve-configs}
\label{sec:twelve-trained}
\label{sec:excluded-cells}
\label{sec:naming}

The four binary components define $2^4 = 16$ combinations (\autoref{fig:ablation-grid}). Four of them enable SimAM without the convolutional stem. Because SimAM re-weights the activations of a convolutional feature map, it has nothing to act on without the stem; these combinations are therefore not weaker models but ill-formed ones, and they are not trained. The remaining twelve configurations are listed in \autoref{tab:twelve-configs}. Configurations 1 to 8 were defined first and configurations 9, 12, 13 and 16 were added to complete the design, so the numbering omits the four excluded combinations (10, 11, 14 and 15). Configuration $n$ is denoted C$n$ in tables and figures.

\begin{figure}[htbp]
	\centering
	\includegraphics[width=0.8\textwidth]{fig4_2_ablation_grid.png}
	\caption[The factorial ablation design]{The $2^4$ factorial design. Each cell is one combination of the four components; twelve are trained and the four that enable SimAM without the convolutional stem are excluded.}
	\label{fig:ablation-grid}
\end{figure}

\begin{table}[htbp]
	\centering
	\small
	\begin{tabular}{@{}l c c c c@{}}
		\hline
		\textbf{Configuration} & \textbf{EVM} & \textbf{SimAM} & \textbf{CNN} & \textbf{TF} \\
		\hline
		C1   & --         & --         & --         & -- \\
		C2   & --         & --         & --         & \checkmark \\
		C3   & --         & --         & \checkmark & -- \\
		C4   & \checkmark & --         & --         & -- \\
		C5   & --         & \checkmark & \checkmark & -- \\
		C6   & --         & \checkmark & \checkmark & \checkmark \\
		C7   & \checkmark & --         & \checkmark & \checkmark \\
		C8   & \checkmark & \checkmark & \checkmark & \checkmark \\
		C9   & --         & --         & \checkmark & \checkmark \\
		C12  & \checkmark & --         & --         & \checkmark \\
		C13  & \checkmark & --         & \checkmark & -- \\
		C16  & \checkmark & \checkmark & \checkmark & -- \\
		\hline
	\end{tabular}
	\caption[The twelve valid ablation configurations]{The twelve valid ablation configurations. EVM: Eulerian video magnification; CNN: convolutional stem; TF: temporal transformer.}
	\label{tab:twelve-configs}
\end{table}

\subsection{Baseline and Control}
\label{sec:baseline-control}

Two configurations serve different purposes and should not be confused. C4 is the baseline of the study: it applies magnification and no learned component, and the contributions of the learned components are ultimately read against it. C1 is not a second baseline but the control for magnification: it is identical to C4 except that magnification is disabled. C8 enables all four components but is otherwise one of twelve configurations and is not privileged in the analysis.

\subsection{Matched Pairs}
\label{sec:matched-pairs}
\label{sec:matched-pairs-table}

A matched pair consists of two configurations that differ in exactly one component. The difference between their scores estimates the contribution of that component with all others held fixed, and the effect reported for a component is the mean over all of its pairs. Magnification and the transformer can be toggled in each of the six valid combinations of the other three components, giving six pairs each. A SimAM pair requires the stem to be present, and a stem pair requires SimAM to be absent, so these two components have four pairs each. Table~\ref{tab:matched-pairs} lists all twenty pairs. Because an effect averaged over four pairs is a noisier estimate than one averaged over six, effects with different pair counts should be compared with care.

\begin{table}[htbp]
	\centering
	\small
	\begin{tabular}{@{}l l r@{}}
		\hline
		\textbf{Component} & \textbf{Pairs (off $\rightarrow$ on)} & \textbf{Count} \\
		\hline
		Transformer & 1$\rightarrow$2, 4$\rightarrow$12, 3$\rightarrow$9, 13$\rightarrow$7, 5$\rightarrow$6, 16$\rightarrow$8 & 6 \\
		EVM         & 1$\rightarrow$4, 3$\rightarrow$13, 9$\rightarrow$7, 6$\rightarrow$8, 5$\rightarrow$16, 2$\rightarrow$12 & 6 \\
		SimAM       & 3$\rightarrow$5, 9$\rightarrow$6, 7$\rightarrow$8, 13$\rightarrow$16 & 4 \\
		CNN stem    & 1$\rightarrow$3, 2$\rightarrow$9, 4$\rightarrow$13, 12$\rightarrow$7 & 4 \\
		\hline
	\end{tabular}
	\caption[Matched pairs by component]{Matched pairs by component, identified by configuration number.}
	\label{tab:matched-pairs}
\end{table}


% =============================================================================
\section{Model Architecture}
\label{sec:architecture}
% =============================================================================

All twelve configurations are branches of a single model in which each component is executed or skipped according to switches fixed at construction; magnification is applied during preprocessing and is not part of the network. Because every configuration shares the same implementation, initialisation and layer order, two configurations that differ in one switch differ in nothing else.

\subsection{Pipeline Overview}
\label{sec:pipeline-overview}
\label{sec:pipeline-in-words}
\label{sec:pipeline-shapes-discussion}

The pipeline consists of two stages that are executed a different number of times (\autoref{fig:pipeline-overview}). \textbf{Stage~1 (preprocessing)} converts each clip into one motion tensor and runs once per clip. \textbf{Stage~2 (training and evaluation)} loads these tensors and fits a model; it runs once for every configuration and every held-out subject. Because Stage~1 is completed before any model is trained, all configurations are compared on identical inputs.

% -----------------------------------------------------------------------------
\begin{figure}[htbp]
	\centering
	\begin{tikzpicture}[
		font=\scriptsize,
		node distance=3.2mm,
		box/.style={
			draw, rounded corners=2pt, align=center,
			text width=34mm, inner sep=3pt,
			fill=black!4},
		switch/.style={box, dashed, fill=black!10},
		term/.style={box, fill=black!14, font=\scriptsize\bfseries},
		arr/.style={-{Stealth[length=1.6mm]}, semithick}
	]

	% ---------------- Stage 1: preprocessing ----------------
	\node[term] (a1) {CASME II dataset\\\normalfont 200\,fps, FACS-coded};
	\node[box, below=of a1] (a2)
		{Clip selection\\\normalfont 156 clips, 25 subjects};
	\node[box, below=of a2] (a3)
		{Frame preparation\\\normalfont greyscale, $224\times224$};
	\node[box, below=of a3] (a4)
		{Uniform sampling\\\normalfont 33 frames per clip};
	\node[switch, below=of a4] (a5)
		{Motion magnification\\\normalfont Eulerian, $\alpha=10$};
	\node[box, below=of a5] (a6)
		{Optical flow\\\normalfont Farneb\"ack, 32 pairs};
	\node[box, below=of a6] (a7)
		{Optical strain\\\normalfont spatial gradients};
	\node[box, below=of a7] (a8)
		{Stack and normalise\\\normalfont $[u,v,\varepsilon]$ to $[0,1]$};
	\node[term, below=of a8] (a9)
		{Saved motion tensor\\\normalfont $(3,32,224,224)$};

	\foreach \i/\j in {a1/a2, a2/a3, a3/a4, a4/a5, a5/a6, a6/a7, a7/a8, a8/a9}
		{\draw[arr] (\i) -- (\j);}

	% ---------------- Stage 2: training and evaluation ----------------
	\node[box, right=20mm of a2] (b1)
		{Load and standardise\\\normalfont flip augmentation};
	\node[switch, below=of b1] (b2)
		{Convolutional stem\\\normalfont three streams};
	\node[switch, below=of b2] (b3)
		{SimAM\\\normalfont parameter-free gate};
	\node[box, below=of b3] (b4)
		{Pool to a sequence\\\normalfont 32 steps, 96 features};
	\node[switch, below=of b4] (b5)
		{Temporal transformer\\\normalfont 4 layers, 8 heads};
	\node[box, below=of b5] (b6)
		{Classifier head\\\normalfont three logits};
	\node[box, below=of b6] (b7)
		{Leave-one-subject-out\\\normalfont 25 folds, counts pooled};
	\node[term, below=of b7] (b8)
		{Pooled macro F1\\\normalfont the primary metric};

	\foreach \i/\j in {b1/b2, b2/b3, b3/b4, b4/b5, b5/b6, b6/b7, b7/b8}
		{\draw[arr] (\i) -- (\j);}

	% ---------------- the hand-off ----------------
	\draw[arr] (a9.east) -- ++(8mm,0) |- (b1.west);

	% ---------------- stage brackets ----------------
	\begin{scope}[on background layer]
		\node[draw=black!45, dashed, rounded corners=3pt, inner sep=4.5mm,
		      fit=(a1)(a9)] (s1) {};
		\node[draw=black!45, dashed, rounded corners=3pt, inner sep=4.5mm,
		      fit=(b1)(b8)] (s2) {};
	\end{scope}
	\node[above=0.5mm of s1, font=\scriptsize\bfseries] {Stage 1: preprocessing};
	\node[above=0.5mm of s2, font=\scriptsize\bfseries] {Stage 2: training and evaluation};

	\end{tikzpicture}
	\caption[The data pipeline, end to end]{The pipeline from dataset to score. Stage~1 runs once per clip; Stage~2 runs once per configuration and fold. Solid boxes are always executed. \textbf{Dashed boxes are the four components switched on and off in the ablation}; each is replaced by a defined operation when disabled.}
	\label{fig:pipeline-overview}
\end{figure}
% -----------------------------------------------------------------------------

In Stage~1, each selected clip is loaded in greyscale, resized to $224\times224$ pixels and sampled at 33 evenly spaced frames between onset and offset. Magnification is applied at this point if enabled. Optical flow is then computed for the 32 adjacent frame pairs, optical strain is derived from its spatial gradients, and the three fields are stacked and normalised into one tensor. The network therefore never receives pixel intensities, only displacement and deformation fields. In Stage~2, each tensor is standardised per channel, optionally flipped for augmentation, and passed through the convolutional stem, the SimAM gate, a pooling step that removes the spatial dimensions, the temporal encoder and the classification head. Table~\ref{tab:pipeline-shapes} lists the representation of the data after each step.

\begin{table}[htbp]
	\centering
	\small
	\begin{tabular}{@{}lll@{}}
		\hline
		\textbf{After this step} & \textbf{The data is} & \textbf{Shape} \\
		\hline
		Clip selection        & a set of clips                & 156 clips, 25 subjects \\
		Frame preparation     & greyscale frames              & $N \times 224 \times 224$ \\
		Uniform sampling      & a fixed-length sequence       & $33 \times 224 \times 224$ \\
		Magnification         & the same, amplified           & $33 \times 224 \times 224$ \\
		Flow and strain       & three motion fields per pair  & $3 \times 32 \times 224 \times 224$ \\
		Pooling to a sequence & a sequence of feature vectors & $32 \times 96$ \\
		Temporal encoding     & one vector per clip           & $96$ \\
		Classifier head       & class scores                  & $3$ \\
		Fold aggregation      & pooled confusion counts       & $3 \times 3$ \\
		\hline
	\end{tabular}
	\caption[What the data is at each pipeline step]{The data representation after each pipeline step, where $N$ is the number of frames between onset and offset. The temporal axis falls from 33 frames to 32 pairs because flow is computed between \emph{adjacent} frames, and the spatial axes are removed at the pooling step rather than at the classifier.}
	\label{tab:pipeline-shapes}
\end{table}

\subsection{Convolutional Stem}
\label{sec:input-streams}
\label{sec:conv-stem}

Each of the three motion channels is processed by its own convolutional stream with unshared weights. Table~\ref{tab:conv-stem} lists the layers of one stream, which use batch normalisation~\cite{ioffe2015}, ReLU activations~\cite{nair2010} and channel-wise dropout~\cite{tompson2015}. No kernel spans the temporal axis (Section~\ref{sec:conv3d}), so the 32 time steps pass through unchanged while height and width are halved once. The outputs of the three streams are concatenated along the channel dimension, giving a feature map of $96\times32\times112\times112$ per clip.

\begin{table}[htbp]
	\centering
	\small
	\begin{tabular}{@{}l l l@{}}
		\hline
		\textbf{Layer} & \textbf{Configuration} & \textbf{Output (per stream)} \\
		\hline
		3D convolution          & 16 filters, kernel $(1,3,3)$, padding $(0,1,1)$, no bias & $16\times32\times224\times224$ \\
		Normalisation and activation & batch normalisation, ReLU, 3D dropout ($p=0.3$) & $16\times32\times224\times224$ \\
		3D convolution          & 32 filters, kernel $(1,3,3)$, padding $(0,1,1)$, no bias & $32\times32\times224\times224$ \\
		Normalisation and activation & batch normalisation, ReLU, 3D dropout ($p=0.3$) & $32\times32\times224\times224$ \\
		3D max pooling          & kernel and stride $(1,2,2)$ & $32\times32\times112\times112$ \\
		\hline
	\end{tabular}
	\caption[Layers of one convolutional stream]{Layers of one convolutional stream. Shapes are given as channels $\times$ time $\times$ height $\times$ width.}
	\label{tab:conv-stem}
\end{table}

\subsection{SimAM}
\label{sec:simam-placement}

When enabled, SimAM~\cite{yang2021} is applied separately to the output of each convolutional stream, before the streams are concatenated. It follows \eqref{eq:simam-stats} and \eqref{eq:simam-gate}, with the mean and variance computed over the time, height and width of each channel (divisor $THW-1$) and $\lambda = 10^{-4}$. It adds no learnable parameters. SimAM is enabled only together with the convolutional stem.

\subsection{Temporal Transformer}
\label{sec:transformer-config}

The concatenated feature map is reduced by adaptive average pooling over the two spatial dimensions to a sequence of 32 time steps, each a 96-dimensional feature vector. Fixed sinusoidal positional encodings~\cite{vaswani2017} are added to this sequence:
\begin{equation}
	\mathrm{PE}(p, 2i) = \sin\!\left(\frac{p}{10000^{2i/96}}\right), \qquad \mathrm{PE}(p, 2i+1) = \cos\!\left(\frac{p}{10000^{2i/96}}\right),
	\label{eq:positional-encoding}
\end{equation}
where $\mathrm{PE}(p, j)$ is the value added to feature $j$ of the token at time step $p \in \{0, \dots, 31\}$, $i$ indexes the pairs of feature dimensions, and 96 is the feature dimension. The encoding is fixed rather than learned, and dropout with probability 0.1 is applied after it is added. The sequence is then processed by a transformer encoder with four layers, eight attention heads, model dimension 96, feed-forward dimension 256, dropout 0.1, GELU activation~\cite{hendrycks2016} and pre-norm layer ordering~\cite{xiong2020}, followed by a final layer normalisation~\cite{ba2016}. The output sequence is averaged over time, giving one 96-dimensional vector per clip.

\subsection{Classification Head and Replacement Operations}
\label{sec:classifier-head}
\label{sec:fallbacks-config}

Every configuration ends in the same classification head: layer normalisation, dropout with probability 0.3, and a linear layer mapping the 96 features to three class logits. A disabled component is replaced by a minimal operation rather than removed, so that the configuration remains a fair control (\autoref{tab:ch2-component-controls}). Without the convolutional stem, each frame of the input tensor is average-pooled to a $4\times4$ grid per channel, and the resulting 48 values are projected to 96 features by a linear layer with 4,704 parameters. Without the transformer, the sequence is averaged over time; this operation has no parameters and no notion of temporal order.

\subsection{Parameter Counts}
\label{sec:parameter-counts}
\label{sec:weight-init}

The convolutional stem has 14,544 parameters across its three streams, the transformer encoder 348,736, the replacement projection 4,704 and the classification head 483; SimAM and magnification add none. Figure~\ref{fig:parameter-counts} shows the resulting totals for the four architectural combinations. In a configuration with all components, the transformer holds approximately 96\,\% of the parameters. All configurations use the same initialisation: Kaiming-normal initialisation for convolutional weights, Xavier-uniform initialisation for linear weights, unit scale and zero shift for normalisation layers, and zero biases.

\begin{figure}[htbp]
	\centering
	\includegraphics[width=0.85\textwidth]{fig4_3_parameter_counts.png}
	\caption[Parameter totals by architectural combination]{Trainable parameters of the four architectural combinations (logarithmic scale), with the contribution of each component.}
	\label{fig:parameter-counts}
\end{figure}


% =============================================================================
\section{Training Procedure}
\label{sec:training-procedure}
% =============================================================================

All twelve configurations are trained with the same procedure, summarised in \autoref{tab:training-settings}; the underlying methods are described in Section~\ref{sec:scarcity-skew}.

\begin{table}[htbp]
	\centering
	\small
	\begin{tabular}{@{}l p{0.62\linewidth}@{}}
		\hline
		\textbf{Setting} & \textbf{Value} \\
		\hline
		Loss            & focal loss~\cite{lin2017} with $\gamma = 2$~\cite{zhaoS2021} and label smoothing~\cite{szegedy2016} of 0.05; no class weights \\
		Sampling        & class-balanced, weight $1/n_c$ per example, with replacement \\
		Augmentation    & random horizontal flip ($p = 0.5$) with negation of the horizontal flow \\
		Optimiser       & AdamW~\cite{loshchilov2019}, learning rate $10^{-4}$, weight decay $10^{-4}$ \\
		Learning rate   & linear warmup over 5 epochs from $10^{-5}$, then cosine annealing~\cite{loshchilov2017} to $10^{-7}$ \\
		Epochs          & 50, without early stopping \\
		Batch size      & 8 \\
		Gradient clipping & global norm 1.0 \\
		Precision       & mixed precision with gradient scaling~\cite{micikevicius2018} \\
		Random seed     & 42, reset before every configuration and fold \\
		\hline
	\end{tabular}
	\caption[Training settings]{Training settings, identical for all configurations and folds.}
	\label{tab:training-settings}
\end{table}

\subsection{Loss and Class Weighting}
\label{sec:loss-config}
\label{sec:class-weighting-rule}
\label{sec:sampler-config}

The loss can apply per-class weights ($\alpha_t$ in \eqref{eq:focal-loss}), but they are enabled only when the class-balanced sampler is inactive. Because the sampler is active in every reported run, the loss carries \textbf{no class weights} (Section~\ref{sec:combining-both}), and the frequency correction comes entirely from the sampler, which gives each training example a weight inversely proportional to the size of its class in the training split of the current fold \eqref{eq:sampler-weight}. The focal term $(1-p_t)^\gamma$ and label smoothing remain active; neither depends on class frequency.

The rule was introduced after two failures in opposite directions, both observed during the development of this study. In an earlier version, the loss carried inverse-frequency class weights and no sampler was used; under this single correction, the proposed configuration stopped predicting the \textbf{Positive} class altogether. A later version added the sampler while keeping the loss weights, so both corrections operated together; the failure then reversed, and configurations stopped predicting the \textbf{Negative} majority class, one of them predicting a single class for every clip. With the rule in place, no configuration in the reported study abandons any class. The only reviewed study that combines two corrections reports that the balanced loss adds only a slight improvement, not that the combination reverses its effect~\cite{xia2020a}. The two failures, however, occurred in implementation versions that differ from the reported run in more than the correction rule, so \textbf{the causal effect of the rule is not isolated}.

\subsection{Optimisation}
\label{sec:optimiser-config}
\label{sec:batch-size}
\label{sec:numerical-safeguards}

Figure~\ref{fig:lr-schedule} shows the learning-rate schedule. The learning rate rises linearly from one tenth of its target value over the first five epochs and then follows a cosine decay to $10^{-7}$; the schedule is updated once per epoch. Every configuration and fold is trained for the full 50 epochs. Gradients are clipped to a global norm of 1.0 after being unscaled from mixed precision, and an update is skipped if any gradient contains a non-finite value.

\begin{figure}[htbp]
	\centering
	\includegraphics[width=0.8\textwidth]{fig4_4_lr_schedule.png}
	\caption[Learning-rate schedule]{Learning rate per epoch: linear warmup over five epochs followed by cosine annealing.}
	\label{fig:lr-schedule}
\end{figure}

\subsection{Reproducibility}
\label{sec:seeding}
\label{sec:no-hyperparameter-search}

The random seed is reset before every configuration and every fold, so the sampling order, weight initialisation and all other stochastic steps start from the same state each time a fold is trained. No hyperparameter search was performed: all training settings are fixed defaults, identical across the twelve configurations and the 25 folds. Every comparison is therefore between architectures trained with one common, untuned recipe, not between architectures each tuned to its own advantage.


% =============================================================================
\section{Evaluation Protocol}
\label{sec:evaluation-protocol}
% =============================================================================

All configurations are evaluated with the same fold structure, checkpoint selection rule and aggregation across folds. The underlying metrics are defined in Section~\ref{sec:evaluation}.

\subsection{Leave-One-Subject-Out Folds}
\label{sec:loso-folds}

Evaluation uses leave-one-subject-out cross-validation~\cite{yan2014,see2019} over the 25 subjects of the working set, giving \textbf{25 folds}. Each fold holds out all clips of one subject and trains on the clips of the remaining 24 subjects. All 25 folds are run for every configuration.

\subsection{Checkpoint Selection}
\label{sec:checkpoint-selection}

Within each fold, the macro F1 on the held-out subject is computed after every epoch, and the weights of the best epoch are restored and scored on the same held-out subject. \textbf{There is no inner validation split}: the held-out clips are used both to select the epoch and to score it. A proper nested procedure would select the epoch on a validation subset of the training subjects and reserve the held-out subject for the final evaluation. The reported scores are therefore optimistically biased by an unknown amount~\cite{varma2006}. The bias is likely to be largest for the smallest folds, where the choice among 50 epochs is made on only a few clips.

\subsection{Aggregation and Primary Metric}
\label{sec:aggregation}
\label{sec:primary-metric}
\label{sec:mean-of-folds-rejected}

The confusion matrices of the 25 folds are summed into one pooled $3\times3$ matrix, from which the per-class precision, recall and F1 are computed. The primary metric is \textbf{pooled macro F1}, the mean of the three pooled per-class F1 scores, which is identical in construction to the Unweighted F1 of MEGC 2019~\cite{see2019}. Pooled accuracy is reported alongside it.

The alternative, averaging the macro F1 of each fold, is rejected because of the fold composition. Of the 25 folds, 10 contain clips of only one class, 8 contain two and 7 contain all three. A class absent from a fold contributes an F1 of zero to that fold, so a single-class fold can score at most $1/3$ and a two-class fold at most $2/3$. The mean-of-folds macro F1 is therefore bounded by
\begin{equation}
	\overline{F1}_{\max} = \frac{10 \cdot \tfrac13 + 8 \cdot \tfrac23 + 7 \cdot 1}{25} = \frac{3.333\ldots + 5.333\ldots + 7}{25} \approx 0.627,
	\label{eq:fold-ceiling}
\end{equation}
where $\overline{F1}_{\max}$ is the highest mean-of-folds macro F1 that any configuration can reach, the numerators weight each fold type by its maximum score, and 25 is the number of folds. No configuration can exceed this ceiling, however accurate it is. Pooling avoids the ceiling, because true positives, false positives and false negatives are accumulated over all 156 clips before any F1 is computed. Fold composition also inflates mean-of-folds accuracy: for the configuration with all four components, it exceeds pooled accuracy by 6.3 percentage points.


% =============================================================================
\section{Research Gaps and Limitations}
\label{sec:gaps-limitations}
% =============================================================================

This section combines, for each component, the research gap left by the reviewed literature (Chapter~\ref{ch:literature}) and the limitations of its implementation in this thesis. The components are grouped into motion modelling, which covers the representation and spatial processing of facial motion, and temporal modelling, which covers the sampling of the sequence and the modelling of its evolution over time. Limitations of the training and evaluation procedure as a whole are listed last.

\subsection{Motion Modelling}
\label{sec:gaps-motion}

\subsubsection{Eulerian Video Magnification}

\textbf{Research gap.} Magnification improves recognition in the reviewed studies, but every evaluation places it in front of an appearance-based representation such as LBP, HOG or pre-trained face features~\cite{liX2018,bai2021,liY2018,liY2021}. No study measures its contribution when the classifier receives an explicit motion field. The magnification factor has an interior optimum that can be found only by search~\cite{liX2018}, and magnification also amplifies unwanted motion such as head movement~\cite{liX2018}.

\textbf{Limitation.} The factor is fixed at $\alpha = 10$, inside the reported optimum, but is not varied. The pass-band, the filter and the order of magnification and sampling depart from the reference formulation (\autoref{tab:ch2-evm-departures}); in particular, the realised pass-band depends on the clip length and is approximately 2.5--12.5~Hz at the median length. No face alignment precedes magnification, so head motion within the band is amplified together with facial movement.

\subsubsection{Optical Flow}

\textbf{Research gap.} Optical flow is the preferred input of the strongest benchmark systems~\cite{see2019}, but the choice of estimator is not settled: Farneb\"ack and TV-L1 are both used~\cite{zhaoS2021,liong2019a}, and no controlled comparison of their recognition performance exists. Whether a single onset-to-apex flow field~\cite{liong2019a,liong2019b} or a sequence of flow fields~\cite{zhaoS2021} is preferable has likewise not been measured.

\textbf{Limitation.} Only the Farneb\"ack estimator~\cite{farneback2003} is used, and flow is computed between consecutive sampled frames rather than relative to the onset. No flow-based fine alignment is performed, although a small misalignment can considerably degrade micro-expression recognition~\cite{xu2017}.

\subsubsection{Optical Strain}

\textbf{Research gap.} The evidence for strain comes from hand-crafted descriptors evaluated under a protocol that is not subject-independent~\cite{liong2016}, and the strain magnitude of STSTNet cannot be determined unambiguously from its published equation~\cite{liong2019b}.

\textbf{Limitation.} The implemented magnitude counts shear once rather than twice, weighting it differently from the Frobenius norm used by Liong et al.~\cite{liong2014a,liong2016}, and no dedicated strain filter is applied~\cite{liong2014a}. Strain is present in all twelve configurations, so its own contribution is not measured.

\subsubsection{Convolutional Spatial Stem}

\textbf{Research gap.} The reviewed studies add modules to a backbone that remains in place~\cite{xia2020b}; none removes the learned spatial stem and passes the motion representation directly to the classifier, so its contribution is unknown. Deeper models and higher input resolutions reduce performance on composite micro-expression data~\cite{xia2020b}.

\textbf{Limitation.} The stem contains only spatial kernels, so the ablation tests a learned spatial stem rather than spatio-temporal convolution~\cite{ji2013,tran2015}. Its input resolution of $224\times224$ exceeds the $100\times100$ threshold reported by Xia et al.~\cite{xia2020b}, so any effect measured for the stem may reflect resolution as much as architecture.

\subsubsection{Parameter-Free Attention}

\textbf{Research gap.} SimAM has been evaluated on natural images with learned appearance features~\cite{yang2021}; its suitability for flow-derived features, and whether its statistics should be computed per frame or per clip, are untested. Its regularisation constant was selected separately for each image dataset~\cite{yang2021}.

\textbf{Limitation.} SimAM statistics are computed over the whole clip rather than per frame, and $\lambda = 10^{-4}$ is not re-tuned. Because SimAM requires the convolutional stem, its effect is measured over only four matched pairs and is conditional on the presence of the stem.

\subsection{Temporal Modelling}
\label{sec:gaps-temporal}

\subsubsection{Temporal Normalisation}

\textbf{Research gap.} Sequence length is usually fixed at ten frames by convention~\cite{liX2013,lu2015}, and the studies that vary it report different optima~\cite{liX2018,ben2021,xu2017}. No reviewed study compares interpolation methods for a motion-based representation, in which synthesised frames would produce synthesised displacements; Xu et al. prefer linear interpolation for motion features but do not compare it with TIM~\cite{xu2017}.

\textbf{Limitation.} Temporal normalisation is implemented as uniform sampling of existing frames rather than the Temporal Interpolation Model~\cite{pfister2011}, so short clips cannot be lengthened and three clips repeat frames. The length of 33 frames is not varied. Sampling reduces the effective frame rate to approximately 100~fps at the median clip length and 52~fps for the longest clip (Section~\ref{sec:three-channel-tensor}).

\subsubsection{Temporal Transformer}

\textbf{Research gap.} The strongest transformer-based results combine a pre-trained spatial transformer with an LSTM aggregator~\cite{zhang2022}, so the benefit of a temporal transformer trained from scratch on one small dataset has not been isolated. Transformers lack the inductive biases of convolution and can perform poorly when trained on insufficient data~\cite{dosovitskiy2021}.

\textbf{Limitation.} The implemented encoder differs from SLSTT in six respects (\autoref{tab:slstt-divergences}). The ablation compares the complete encoder, including positional encoding, feed-forward layers and normalisation, with mean pooling; it does not isolate self-attention or compare it with another learned temporal model.

\subsection{Training and Evaluation}
\label{sec:gaps-training}

\textbf{Research gap.} How imbalance corrections interact is examined by only one reviewed study~\cite{xia2020a}, the focusing parameter of focal loss has not been validated for this task~\cite{zhaoS2021}, and label smoothing~\cite{szegedy2016} has not been examined for micro-expression recognition.

\textbf{Limitation.} The imbalance treatment, $\gamma$ and label smoothing are identical in all configurations, so their effects are not measured. All results come from one dataset, one random seed and a checkpoint selection without an inner validation split (Sections~\ref{sec:design-limitation} and~\ref{sec:checkpoint-selection}), and no result extends to other datasets, populations or recording conditions (Section~\ref{sec:sibling-corpora}).


% =============================================================================
\section{Summary}
\label{sec:ch4-conclusion}
% =============================================================================

This chapter described the factorial ablation of four pipeline components on CASME~II. It presented the 156-clip working set and its 25 leave-one-subject-out folds, the twelve valid configurations and the twenty matched pairs used to compare them, and the shared model architecture, training procedure and evaluation protocol, with pooled macro F1 as the primary metric. It closed with the research gaps and limitations of each motion-modelling and temporal-modelling component.

Chapter~\ref{ch:results} reports the results of the ablation: the performance of the twelve configurations and the contribution of each component, measured over its matched pairs.
. This file
% is the single source for Chapter 4; chapter4.tex and the chapter4_short_half
% files are superseded and share labels with it, so never include them.
%
% Rewritten 23 September 2026:
%   * Subsections consolidated; merged subsections keep every earlier label
%     (stacked \label commands), so references from other chapters resolve.
%   * Removed: "What is stored" and "Earlier evaluations" (labels
%     sec:what-is-stored, sec:excluded-evaluations, tab:five-evaluations),
%     all VRAM figures, the Computational Cost section (removed at the
%     author's request), and every code file, script and class name.
%   * The pipeline overview (formerly its own section) is now Section 4.4.1.
%   * New Section 4.8 combines the research gap and the limitations of each
%     component, grouped into motion modelling and temporal modelling.
%   * Figures fig4_1 ... fig4_4 are generated by tools/thesis_ch4_figures.py
%     from the values stated in this chapter.
%
% Requires in the preamble: amsmath, amssymb (\checkmark), graphicx, tikz with
%   \usetikzlibrary{arrows.meta, positioning, fit, backgrounds}
% =============================================================================

\chapter{Methodology}
\label{ch:methodology}

This chapter describes the study as it was carried out. Section~\ref{sec:research-approach} presents the experimental design, Section~\ref{sec:corpus-preparation} the dataset and its preparation, and Section~\ref{sec:ablation-matrix} the twelve configurations and the matched pairs used to compare them. Section~\ref{sec:architecture} describes the model architecture, Section~\ref{sec:training-procedure} the training procedure, and Section~\ref{sec:evaluation-protocol} the evaluation protocol. Section~\ref{sec:gaps-limitations} summarises the research gaps and limitations of each component, grouped into motion modelling and temporal modelling.


% =============================================================================
\section{Research Approach and Experimental Design}
\label{sec:research-approach}
% =============================================================================

\subsection{Factorial Ablation Design}
\label{sec:study-form}

This thesis does not propose and evaluate a single architecture. The study is designed as a factorial ablation: four components of a micro-expression recognition pipeline are each switched on or off, every valid combination is trained and evaluated under an identical protocol, and the contribution of each component is estimated from pairs of configurations that differ in that component alone. The object of study is therefore the set of configurations, and the question is not how well the pipeline performs but which of its stages are responsible for that performance, and by how much.

This design addresses the research gap of Section~\ref{sec:research-gap}. In the reviewed literature, components are varied one at a time within a single proposed system, for example the magnification factor~\cite{liX2018}, individual modules of a recurrent convolutional network~\cite{xia2020b} or the temporal aggregator of a transformer~\cite{zhang2022}. No reviewed study varies several components jointly under one protocol with all other factors held identical.

\subsection{Ablated Components}
\label{sec:four-variables}

The design space is defined by four components, each of which is either fully present or absent:

\begin{itemize}
	\item \textbf{Eulerian video magnification (EVM)}~\cite{wu2012}, applied during preprocessing before the motion representation is computed (Section~\ref{sec:motion-magnification}).
	\item \textbf{A convolutional spatial stem}, with one stream per motion channel (Section~\ref{sec:conv3d}).
	\item \textbf{SimAM}~\cite{yang2021}, a parameter-free attention module within the convolutional streams (Section~\ref{sec:simam}).
	\item \textbf{A temporal transformer encoder}~\cite{vaswani2017}, applied to the sequence of frame features (Section~\ref{sec:self-attention}).
\end{itemize}

The first three act on the motion representation and are grouped as \textit{motion modelling}; the transformer models the evolution of that representation over time and constitutes \textit{temporal modelling}. A disabled component is replaced by a minimal operation (Section~\ref{sec:fallbacks-config}) rather than by an alternative component of comparable capacity.

\subsection{Design Limitations}
\label{sec:design-limitation}
\label{sec:what-is-stored}

Every configuration is trained once, under a single fixed random seed (42). No configuration is retrained with other seeds, so no variance across seeds is estimated. Every difference between two configurations, including the matched-pair effects of Section~\ref{sec:matched-pairs}, is therefore a difference between two point estimates without a confidence interval, and it could change under a different seed. In addition, per-clip predictions were not retained; only aggregate metrics and confusion matrices were stored. A paired significance test such as McNemar's test~\cite{mcnemar1947}, which requires the prediction for each clip, therefore cannot be computed for any pair of configurations.


% =============================================================================
\section{Dataset and Preparation}
\label{sec:corpus-preparation}
% =============================================================================

\subsection{Data Selection}
\label{sec:corpus-used}
\label{sec:255-to-156}

All experiments use CASME~II~\cite{yan2014} (Section~\ref{sec:corpus}); no other dataset contributes to training or evaluation, so every result is a within-dataset result. The released coding contains 255 clips coded as micro-expressions with extractable frames (Section~\ref{sec:clip-selection}). The \texttt{Others} category (99 clips) is excluded before training. It is a residual category for clips that do not fit the other labels, and a class defined by exclusion rather than content cannot provide a meaningful training signal. Figure~\ref{fig:class-distribution} shows the composition of the 255 clips and the resulting working set of 156 clips.

\begin{figure}[htbp]
	\centering
	\includegraphics[width=0.85\textwidth]{fig4_1_class_distribution.png}
	\caption[Composition of the CASME II working set]{Composition of the 255 CASME~II micro-expression clips by emotion label, grouped into the three classes used in this thesis. The \texttt{Others} category is excluded, leaving a working set of 156 clips. Counts are taken from the released coding (Section~\ref{sec:label-taxonomy}).}
	\label{fig:class-distribution}
\end{figure}

\subsection{Class Grouping}
\label{sec:three-class-grouping}

Disgust, repression, sadness and fear are grouped into Negative, happiness forms Positive and surprise forms Surprise, following the valence-based convention of SMIC and MEGC 2019~\cite{liX2013,see2019} (Section~\ref{sec:label-taxonomy}). The resulting class distribution is 99 : 32 : 25, approximately 4 : 1.3 : 1. All results are reported for these three classes.

\subsection{Subjects and Folds}
\label{sec:subjects-loso}

The 255 clips come from 26 subjects. Subject 18 contributes only \texttt{Others} clips and is removed together with that category, so the working set contains 25 subjects, and leave-one-subject-out cross-validation produces \textbf{25 folds} rather than 26. Because subjects contribute unequal numbers of clips, the folds differ considerably in size and class composition (\autoref{fig:fold-composition}).

\begin{figure}[tbp]
	\centering
	\includegraphics[width=\textwidth]{fig3_2_fold_composition.png}
	\caption[Class composition of the 25 leave-one-subject-out folds]{Class composition of the 25 leave-one-subject-out folds. Top: clips held out in each fold, by class, with the number of classes present above each bar. Bottom: the number of folds containing one, two or three classes, and the per-fold macro-F1 ceiling this implies.}
	\label{fig:fold-composition}
\end{figure}

\subsection{Input Preparation}
\label{sec:network-preparation}

Each clip is converted into a motion tensor of shape $(3, 32, 224, 224)$, whose three channels hold horizontal optical flow, vertical optical flow and optical strain magnitude over 32 frame pairs (Sections~\ref{sec:recording-to-input} and~\ref{sec:video-to-motion}). Table~\ref{tab:preprocessing-settings} lists the settings used. None of them is varied by the ablation, so every configuration receives inputs of the same shape.

\begin{table}[htbp]
	\centering
	\small
	\begin{tabular}{@{}l l@{}}
		\hline
		\textbf{Setting} & \textbf{Value} \\
		\hline
		Frames sampled per clip        & 33, uniformly between onset and offset (32 flow pairs) \\
		Frame size                     & $224\times224$ pixels, greyscale \\
		Magnification factor $\alpha$  & 10 \\
		Magnification band             & 5--25 Hz at a nominal 200 fps \\
		Laplacian pyramid levels       & 4 \\
		Magnified output               & clipped to $[0,255]$ and converted to 8-bit \\
		Optical flow estimator         & Farneb\"ack~\cite{farneback2003} \\
		Channel normalisation          & min--max to $[0,1]$ per clip \\
		\hline
	\end{tabular}
	\caption[Preprocessing settings]{Preprocessing settings, fixed across all configurations. Section~\ref{sec:evm-departures} lists how the magnification settings depart from the reference formulation.}
	\label{tab:preprocessing-settings}
\end{table}


% =============================================================================
\section{Ablation Configurations}
\label{sec:ablation-matrix}
% =============================================================================

\subsection{Valid Configurations}
\label{sec:twelve-configs}
\label{sec:twelve-trained}
\label{sec:excluded-cells}
\label{sec:naming}

The four binary components define $2^4 = 16$ combinations (\autoref{fig:ablation-grid}). Four of them enable SimAM without the convolutional stem. Because SimAM re-weights the activations of a convolutional feature map, it has nothing to act on without the stem; these combinations are therefore not weaker models but ill-formed ones, and they are not trained. The remaining twelve configurations are listed in \autoref{tab:twelve-configs}. Configurations 1 to 8 were defined first and configurations 9, 12, 13 and 16 were added to complete the design, so the numbering omits the four excluded combinations (10, 11, 14 and 15). Configuration $n$ is denoted C$n$ in tables and figures.

\begin{figure}[htbp]
	\centering
	\includegraphics[width=0.8\textwidth]{fig4_2_ablation_grid.png}
	\caption[The factorial ablation design]{The $2^4$ factorial design. Each cell is one combination of the four components; twelve are trained and the four that enable SimAM without the convolutional stem are excluded.}
	\label{fig:ablation-grid}
\end{figure}

\begin{table}[htbp]
	\centering
	\small
	\begin{tabular}{@{}l c c c c@{}}
		\hline
		\textbf{Configuration} & \textbf{EVM} & \textbf{SimAM} & \textbf{CNN} & \textbf{TF} \\
		\hline
		C1   & --         & --         & --         & -- \\
		C2   & --         & --         & --         & \checkmark \\
		C3   & --         & --         & \checkmark & -- \\
		C4   & \checkmark & --         & --         & -- \\
		C5   & --         & \checkmark & \checkmark & -- \\
		C6   & --         & \checkmark & \checkmark & \checkmark \\
		C7   & \checkmark & --         & \checkmark & \checkmark \\
		C8   & \checkmark & \checkmark & \checkmark & \checkmark \\
		C9   & --         & --         & \checkmark & \checkmark \\
		C12  & \checkmark & --         & --         & \checkmark \\
		C13  & \checkmark & --         & \checkmark & -- \\
		C16  & \checkmark & \checkmark & \checkmark & -- \\
		\hline
	\end{tabular}
	\caption[The twelve valid ablation configurations]{The twelve valid ablation configurations. EVM: Eulerian video magnification; CNN: convolutional stem; TF: temporal transformer.}
	\label{tab:twelve-configs}
\end{table}

\subsection{Baseline and Control}
\label{sec:baseline-control}

Two configurations serve different purposes and should not be confused. C4 is the baseline of the study: it applies magnification and no learned component, and the contributions of the learned components are ultimately read against it. C1 is not a second baseline but the control for magnification: it is identical to C4 except that magnification is disabled. C8 enables all four components but is otherwise one of twelve configurations and is not privileged in the analysis.

\subsection{Matched Pairs}
\label{sec:matched-pairs}
\label{sec:matched-pairs-table}

A matched pair consists of two configurations that differ in exactly one component. The difference between their scores estimates the contribution of that component with all others held fixed, and the effect reported for a component is the mean over all of its pairs. Magnification and the transformer can be toggled in each of the six valid combinations of the other three components, giving six pairs each. A SimAM pair requires the stem to be present, and a stem pair requires SimAM to be absent, so these two components have four pairs each. Table~\ref{tab:matched-pairs} lists all twenty pairs. Because an effect averaged over four pairs is a noisier estimate than one averaged over six, effects with different pair counts should be compared with care.

\begin{table}[htbp]
	\centering
	\small
	\begin{tabular}{@{}l l r@{}}
		\hline
		\textbf{Component} & \textbf{Pairs (off $\rightarrow$ on)} & \textbf{Count} \\
		\hline
		Transformer & 1$\rightarrow$2, 4$\rightarrow$12, 3$\rightarrow$9, 13$\rightarrow$7, 5$\rightarrow$6, 16$\rightarrow$8 & 6 \\
		EVM         & 1$\rightarrow$4, 3$\rightarrow$13, 9$\rightarrow$7, 6$\rightarrow$8, 5$\rightarrow$16, 2$\rightarrow$12 & 6 \\
		SimAM       & 3$\rightarrow$5, 9$\rightarrow$6, 7$\rightarrow$8, 13$\rightarrow$16 & 4 \\
		CNN stem    & 1$\rightarrow$3, 2$\rightarrow$9, 4$\rightarrow$13, 12$\rightarrow$7 & 4 \\
		\hline
	\end{tabular}
	\caption[Matched pairs by component]{Matched pairs by component, identified by configuration number.}
	\label{tab:matched-pairs}
\end{table}


% =============================================================================
\section{Model Architecture}
\label{sec:architecture}
% =============================================================================

All twelve configurations are branches of a single model in which each component is executed or skipped according to switches fixed at construction; magnification is applied during preprocessing and is not part of the network. Because every configuration shares the same implementation, initialisation and layer order, two configurations that differ in one switch differ in nothing else.

\subsection{Pipeline Overview}
\label{sec:pipeline-overview}
\label{sec:pipeline-in-words}
\label{sec:pipeline-shapes-discussion}

The pipeline consists of two stages that are executed a different number of times (\autoref{fig:pipeline-overview}). \textbf{Stage~1 (preprocessing)} converts each clip into one motion tensor and runs once per clip. \textbf{Stage~2 (training and evaluation)} loads these tensors and fits a model; it runs once for every configuration and every held-out subject. Because Stage~1 is completed before any model is trained, all configurations are compared on identical inputs.

% -----------------------------------------------------------------------------
\begin{figure}[htbp]
	\centering
	\begin{tikzpicture}[
		font=\scriptsize,
		node distance=3.2mm,
		box/.style={
			draw, rounded corners=2pt, align=center,
			text width=34mm, inner sep=3pt,
			fill=black!4},
		switch/.style={box, dashed, fill=black!10},
		term/.style={box, fill=black!14, font=\scriptsize\bfseries},
		arr/.style={-{Stealth[length=1.6mm]}, semithick}
	]

	% ---------------- Stage 1: preprocessing ----------------
	\node[term] (a1) {CASME II dataset\\\normalfont 200\,fps, FACS-coded};
	\node[box, below=of a1] (a2)
		{Clip selection\\\normalfont 156 clips, 25 subjects};
	\node[box, below=of a2] (a3)
		{Frame preparation\\\normalfont greyscale, $224\times224$};
	\node[box, below=of a3] (a4)
		{Uniform sampling\\\normalfont 33 frames per clip};
	\node[switch, below=of a4] (a5)
		{Motion magnification\\\normalfont Eulerian, $\alpha=10$};
	\node[box, below=of a5] (a6)
		{Optical flow\\\normalfont Farneb\"ack, 32 pairs};
	\node[box, below=of a6] (a7)
		{Optical strain\\\normalfont spatial gradients};
	\node[box, below=of a7] (a8)
		{Stack and normalise\\\normalfont $[u,v,\varepsilon]$ to $[0,1]$};
	\node[term, below=of a8] (a9)
		{Saved motion tensor\\\normalfont $(3,32,224,224)$};

	\foreach \i/\j in {a1/a2, a2/a3, a3/a4, a4/a5, a5/a6, a6/a7, a7/a8, a8/a9}
		{\draw[arr] (\i) -- (\j);}

	% ---------------- Stage 2: training and evaluation ----------------
	\node[box, right=20mm of a2] (b1)
		{Load and standardise\\\normalfont flip augmentation};
	\node[switch, below=of b1] (b2)
		{Convolutional stem\\\normalfont three streams};
	\node[switch, below=of b2] (b3)
		{SimAM\\\normalfont parameter-free gate};
	\node[box, below=of b3] (b4)
		{Pool to a sequence\\\normalfont 32 steps, 96 features};
	\node[switch, below=of b4] (b5)
		{Temporal transformer\\\normalfont 4 layers, 8 heads};
	\node[box, below=of b5] (b6)
		{Classifier head\\\normalfont three logits};
	\node[box, below=of b6] (b7)
		{Leave-one-subject-out\\\normalfont 25 folds, counts pooled};
	\node[term, below=of b7] (b8)
		{Pooled macro F1\\\normalfont the primary metric};

	\foreach \i/\j in {b1/b2, b2/b3, b3/b4, b4/b5, b5/b6, b6/b7, b7/b8}
		{\draw[arr] (\i) -- (\j);}

	% ---------------- the hand-off ----------------
	\draw[arr] (a9.east) -- ++(8mm,0) |- (b1.west);

	% ---------------- stage brackets ----------------
	\begin{scope}[on background layer]
		\node[draw=black!45, dashed, rounded corners=3pt, inner sep=4.5mm,
		      fit=(a1)(a9)] (s1) {};
		\node[draw=black!45, dashed, rounded corners=3pt, inner sep=4.5mm,
		      fit=(b1)(b8)] (s2) {};
	\end{scope}
	\node[above=0.5mm of s1, font=\scriptsize\bfseries] {Stage 1: preprocessing};
	\node[above=0.5mm of s2, font=\scriptsize\bfseries] {Stage 2: training and evaluation};

	\end{tikzpicture}
	\caption[The data pipeline, end to end]{The pipeline from dataset to score. Stage~1 runs once per clip; Stage~2 runs once per configuration and fold. Solid boxes are always executed. \textbf{Dashed boxes are the four components switched on and off in the ablation}; each is replaced by a defined operation when disabled.}
	\label{fig:pipeline-overview}
\end{figure}
% -----------------------------------------------------------------------------

In Stage~1, each selected clip is loaded in greyscale, resized to $224\times224$ pixels and sampled at 33 evenly spaced frames between onset and offset. Magnification is applied at this point if enabled. Optical flow is then computed for the 32 adjacent frame pairs, optical strain is derived from its spatial gradients, and the three fields are stacked and normalised into one tensor. The network therefore never receives pixel intensities, only displacement and deformation fields. In Stage~2, each tensor is standardised per channel, optionally flipped for augmentation, and passed through the convolutional stem, the SimAM gate, a pooling step that removes the spatial dimensions, the temporal encoder and the classification head. Table~\ref{tab:pipeline-shapes} lists the representation of the data after each step.

\begin{table}[htbp]
	\centering
	\small
	\begin{tabular}{@{}lll@{}}
		\hline
		\textbf{After this step} & \textbf{The data is} & \textbf{Shape} \\
		\hline
		Clip selection        & a set of clips                & 156 clips, 25 subjects \\
		Frame preparation     & greyscale frames              & $N \times 224 \times 224$ \\
		Uniform sampling      & a fixed-length sequence       & $33 \times 224 \times 224$ \\
		Magnification         & the same, amplified           & $33 \times 224 \times 224$ \\
		Flow and strain       & three motion fields per pair  & $3 \times 32 \times 224 \times 224$ \\
		Pooling to a sequence & a sequence of feature vectors & $32 \times 96$ \\
		Temporal encoding     & one vector per clip           & $96$ \\
		Classifier head       & class scores                  & $3$ \\
		Fold aggregation      & pooled confusion counts       & $3 \times 3$ \\
		\hline
	\end{tabular}
	\caption[What the data is at each pipeline step]{The data representation after each pipeline step, where $N$ is the number of frames between onset and offset. The temporal axis falls from 33 frames to 32 pairs because flow is computed between \emph{adjacent} frames, and the spatial axes are removed at the pooling step rather than at the classifier.}
	\label{tab:pipeline-shapes}
\end{table}

\subsection{Convolutional Stem}
\label{sec:input-streams}
\label{sec:conv-stem}

Each of the three motion channels is processed by its own convolutional stream with unshared weights. Table~\ref{tab:conv-stem} lists the layers of one stream, which use batch normalisation~\cite{ioffe2015}, ReLU activations~\cite{nair2010} and channel-wise dropout~\cite{tompson2015}. No kernel spans the temporal axis (Section~\ref{sec:conv3d}), so the 32 time steps pass through unchanged while height and width are halved once. The outputs of the three streams are concatenated along the channel dimension, giving a feature map of $96\times32\times112\times112$ per clip.

\begin{table}[htbp]
	\centering
	\small
	\begin{tabular}{@{}l l l@{}}
		\hline
		\textbf{Layer} & \textbf{Configuration} & \textbf{Output (per stream)} \\
		\hline
		3D convolution          & 16 filters, kernel $(1,3,3)$, padding $(0,1,1)$, no bias & $16\times32\times224\times224$ \\
		Normalisation and activation & batch normalisation, ReLU, 3D dropout ($p=0.3$) & $16\times32\times224\times224$ \\
		3D convolution          & 32 filters, kernel $(1,3,3)$, padding $(0,1,1)$, no bias & $32\times32\times224\times224$ \\
		Normalisation and activation & batch normalisation, ReLU, 3D dropout ($p=0.3$) & $32\times32\times224\times224$ \\
		3D max pooling          & kernel and stride $(1,2,2)$ & $32\times32\times112\times112$ \\
		\hline
	\end{tabular}
	\caption[Layers of one convolutional stream]{Layers of one convolutional stream. Shapes are given as channels $\times$ time $\times$ height $\times$ width.}
	\label{tab:conv-stem}
\end{table}

\subsection{SimAM}
\label{sec:simam-placement}

When enabled, SimAM~\cite{yang2021} is applied separately to the output of each convolutional stream, before the streams are concatenated. It follows \eqref{eq:simam-stats} and \eqref{eq:simam-gate}, with the mean and variance computed over the time, height and width of each channel (divisor $THW-1$) and $\lambda = 10^{-4}$. It adds no learnable parameters. SimAM is enabled only together with the convolutional stem.

\subsection{Temporal Transformer}
\label{sec:transformer-config}

The concatenated feature map is reduced by adaptive average pooling over the two spatial dimensions to a sequence of 32 time steps, each a 96-dimensional feature vector. Fixed sinusoidal positional encodings~\cite{vaswani2017} are added to this sequence:
\begin{equation}
	\mathrm{PE}(p, 2i) = \sin\!\left(\frac{p}{10000^{2i/96}}\right), \qquad \mathrm{PE}(p, 2i+1) = \cos\!\left(\frac{p}{10000^{2i/96}}\right),
	\label{eq:positional-encoding}
\end{equation}
where $\mathrm{PE}(p, j)$ is the value added to feature $j$ of the token at time step $p \in \{0, \dots, 31\}$, $i$ indexes the pairs of feature dimensions, and 96 is the feature dimension. The encoding is fixed rather than learned, and dropout with probability 0.1 is applied after it is added. The sequence is then processed by a transformer encoder with four layers, eight attention heads, model dimension 96, feed-forward dimension 256, dropout 0.1, GELU activation~\cite{hendrycks2016} and pre-norm layer ordering~\cite{xiong2020}, followed by a final layer normalisation~\cite{ba2016}. The output sequence is averaged over time, giving one 96-dimensional vector per clip.

\subsection{Classification Head and Replacement Operations}
\label{sec:classifier-head}
\label{sec:fallbacks-config}

Every configuration ends in the same classification head: layer normalisation, dropout with probability 0.3, and a linear layer mapping the 96 features to three class logits. A disabled component is replaced by a minimal operation rather than removed, so that the configuration remains a fair control (\autoref{tab:ch2-component-controls}). Without the convolutional stem, each frame of the input tensor is average-pooled to a $4\times4$ grid per channel, and the resulting 48 values are projected to 96 features by a linear layer with 4,704 parameters. Without the transformer, the sequence is averaged over time; this operation has no parameters and no notion of temporal order.

\subsection{Parameter Counts}
\label{sec:parameter-counts}
\label{sec:weight-init}

The convolutional stem has 14,544 parameters across its three streams, the transformer encoder 348,736, the replacement projection 4,704 and the classification head 483; SimAM and magnification add none. Figure~\ref{fig:parameter-counts} shows the resulting totals for the four architectural combinations. In a configuration with all components, the transformer holds approximately 96\,\% of the parameters. All configurations use the same initialisation: Kaiming-normal initialisation for convolutional weights, Xavier-uniform initialisation for linear weights, unit scale and zero shift for normalisation layers, and zero biases.

\begin{figure}[htbp]
	\centering
	\includegraphics[width=0.85\textwidth]{fig4_3_parameter_counts.png}
	\caption[Parameter totals by architectural combination]{Trainable parameters of the four architectural combinations (logarithmic scale), with the contribution of each component.}
	\label{fig:parameter-counts}
\end{figure}


% =============================================================================
\section{Training Procedure}
\label{sec:training-procedure}
% =============================================================================

All twelve configurations are trained with the same procedure, summarised in \autoref{tab:training-settings}; the underlying methods are described in Section~\ref{sec:scarcity-skew}.

\begin{table}[htbp]
	\centering
	\small
	\begin{tabular}{@{}l p{0.62\linewidth}@{}}
		\hline
		\textbf{Setting} & \textbf{Value} \\
		\hline
		Loss            & focal loss~\cite{lin2017} with $\gamma = 2$~\cite{zhaoS2021} and label smoothing~\cite{szegedy2016} of 0.05; no class weights \\
		Sampling        & class-balanced, weight $1/n_c$ per example, with replacement \\
		Augmentation    & random horizontal flip ($p = 0.5$) with negation of the horizontal flow \\
		Optimiser       & AdamW~\cite{loshchilov2019}, learning rate $10^{-4}$, weight decay $10^{-4}$ \\
		Learning rate   & linear warmup over 5 epochs from $10^{-5}$, then cosine annealing~\cite{loshchilov2017} to $10^{-7}$ \\
		Epochs          & 50, without early stopping \\
		Batch size      & 8 \\
		Gradient clipping & global norm 1.0 \\
		Precision       & mixed precision with gradient scaling~\cite{micikevicius2018} \\
		Random seed     & 42, reset before every configuration and fold \\
		\hline
	\end{tabular}
	\caption[Training settings]{Training settings, identical for all configurations and folds.}
	\label{tab:training-settings}
\end{table}

\subsection{Loss and Class Weighting}
\label{sec:loss-config}
\label{sec:class-weighting-rule}
\label{sec:sampler-config}

The loss can apply per-class weights ($\alpha_t$ in \eqref{eq:focal-loss}), but they are enabled only when the class-balanced sampler is inactive. Because the sampler is active in every reported run, the loss carries \textbf{no class weights} (Section~\ref{sec:combining-both}), and the frequency correction comes entirely from the sampler, which gives each training example a weight inversely proportional to the size of its class in the training split of the current fold \eqref{eq:sampler-weight}. The focal term $(1-p_t)^\gamma$ and label smoothing remain active; neither depends on class frequency.

The rule was introduced after two failures in opposite directions, both observed during the development of this study. In an earlier version, the loss carried inverse-frequency class weights and no sampler was used; under this single correction, the proposed configuration stopped predicting the \textbf{Positive} class altogether. A later version added the sampler while keeping the loss weights, so both corrections operated together; the failure then reversed, and configurations stopped predicting the \textbf{Negative} majority class, one of them predicting a single class for every clip. With the rule in place, no configuration in the reported study abandons any class. The only reviewed study that combines two corrections reports that the balanced loss adds only a slight improvement, not that the combination reverses its effect~\cite{xia2020a}. The two failures, however, occurred in implementation versions that differ from the reported run in more than the correction rule, so \textbf{the causal effect of the rule is not isolated}.

\subsection{Optimisation}
\label{sec:optimiser-config}
\label{sec:batch-size}
\label{sec:numerical-safeguards}

Figure~\ref{fig:lr-schedule} shows the learning-rate schedule. The learning rate rises linearly from one tenth of its target value over the first five epochs and then follows a cosine decay to $10^{-7}$; the schedule is updated once per epoch. Every configuration and fold is trained for the full 50 epochs. Gradients are clipped to a global norm of 1.0 after being unscaled from mixed precision, and an update is skipped if any gradient contains a non-finite value.

\begin{figure}[htbp]
	\centering
	\includegraphics[width=0.8\textwidth]{fig4_4_lr_schedule.png}
	\caption[Learning-rate schedule]{Learning rate per epoch: linear warmup over five epochs followed by cosine annealing.}
	\label{fig:lr-schedule}
\end{figure}

\subsection{Reproducibility}
\label{sec:seeding}
\label{sec:no-hyperparameter-search}

The random seed is reset before every configuration and every fold, so the sampling order, weight initialisation and all other stochastic steps start from the same state each time a fold is trained. No hyperparameter search was performed: all training settings are fixed defaults, identical across the twelve configurations and the 25 folds. Every comparison is therefore between architectures trained with one common, untuned recipe, not between architectures each tuned to its own advantage.


% =============================================================================
\section{Evaluation Protocol}
\label{sec:evaluation-protocol}
% =============================================================================

All configurations are evaluated with the same fold structure, checkpoint selection rule and aggregation across folds. The underlying metrics are defined in Section~\ref{sec:evaluation}.

\subsection{Leave-One-Subject-Out Folds}
\label{sec:loso-folds}

Evaluation uses leave-one-subject-out cross-validation~\cite{yan2014,see2019} over the 25 subjects of the working set, giving \textbf{25 folds}. Each fold holds out all clips of one subject and trains on the clips of the remaining 24 subjects. All 25 folds are run for every configuration.

\subsection{Checkpoint Selection}
\label{sec:checkpoint-selection}

Within each fold, the macro F1 on the held-out subject is computed after every epoch, and the weights of the best epoch are restored and scored on the same held-out subject. \textbf{There is no inner validation split}: the held-out clips are used both to select the epoch and to score it. A proper nested procedure would select the epoch on a validation subset of the training subjects and reserve the held-out subject for the final evaluation. The reported scores are therefore optimistically biased by an unknown amount~\cite{varma2006}. The bias is likely to be largest for the smallest folds, where the choice among 50 epochs is made on only a few clips.

\subsection{Aggregation and Primary Metric}
\label{sec:aggregation}
\label{sec:primary-metric}
\label{sec:mean-of-folds-rejected}

The confusion matrices of the 25 folds are summed into one pooled $3\times3$ matrix, from which the per-class precision, recall and F1 are computed. The primary metric is \textbf{pooled macro F1}, the mean of the three pooled per-class F1 scores, which is identical in construction to the Unweighted F1 of MEGC 2019~\cite{see2019}. Pooled accuracy is reported alongside it.

The alternative, averaging the macro F1 of each fold, is rejected because of the fold composition. Of the 25 folds, 10 contain clips of only one class, 8 contain two and 7 contain all three. A class absent from a fold contributes an F1 of zero to that fold, so a single-class fold can score at most $1/3$ and a two-class fold at most $2/3$. The mean-of-folds macro F1 is therefore bounded by
\begin{equation}
	\overline{F1}_{\max} = \frac{10 \cdot \tfrac13 + 8 \cdot \tfrac23 + 7 \cdot 1}{25} = \frac{3.333\ldots + 5.333\ldots + 7}{25} \approx 0.627,
	\label{eq:fold-ceiling}
\end{equation}
where $\overline{F1}_{\max}$ is the highest mean-of-folds macro F1 that any configuration can reach, the numerators weight each fold type by its maximum score, and 25 is the number of folds. No configuration can exceed this ceiling, however accurate it is. Pooling avoids the ceiling, because true positives, false positives and false negatives are accumulated over all 156 clips before any F1 is computed. Fold composition also inflates mean-of-folds accuracy: for the configuration with all four components, it exceeds pooled accuracy by 6.3 percentage points.


% =============================================================================
\section{Research Gaps and Limitations}
\label{sec:gaps-limitations}
% =============================================================================

This section combines, for each component, the research gap left by the reviewed literature (Chapter~\ref{ch:literature}) and the limitations of its implementation in this thesis. The components are grouped into motion modelling, which covers the representation and spatial processing of facial motion, and temporal modelling, which covers the sampling of the sequence and the modelling of its evolution over time. Limitations of the training and evaluation procedure as a whole are listed last.

\subsection{Motion Modelling}
\label{sec:gaps-motion}

\subsubsection{Eulerian Video Magnification}

\textbf{Research gap.} Magnification improves recognition in the reviewed studies, but every evaluation places it in front of an appearance-based representation such as LBP, HOG or pre-trained face features~\cite{liX2018,bai2021,liY2018,liY2021}. No study measures its contribution when the classifier receives an explicit motion field. The magnification factor has an interior optimum that can be found only by search~\cite{liX2018}, and magnification also amplifies unwanted motion such as head movement~\cite{liX2018}.

\textbf{Limitation.} The factor is fixed at $\alpha = 10$, inside the reported optimum, but is not varied. The pass-band, the filter and the order of magnification and sampling depart from the reference formulation (\autoref{tab:ch2-evm-departures}); in particular, the realised pass-band depends on the clip length and is approximately 2.5--12.5~Hz at the median length. No face alignment precedes magnification, so head motion within the band is amplified together with facial movement.

\subsubsection{Optical Flow}

\textbf{Research gap.} Optical flow is the preferred input of the strongest benchmark systems~\cite{see2019}, but the choice of estimator is not settled: Farneb\"ack and TV-L1 are both used~\cite{zhaoS2021,liong2019a}, and no controlled comparison of their recognition performance exists. Whether a single onset-to-apex flow field~\cite{liong2019a,liong2019b} or a sequence of flow fields~\cite{zhaoS2021} is preferable has likewise not been measured.

\textbf{Limitation.} Only the Farneb\"ack estimator~\cite{farneback2003} is used, and flow is computed between consecutive sampled frames rather than relative to the onset. No flow-based fine alignment is performed, although a small misalignment can considerably degrade micro-expression recognition~\cite{xu2017}.

\subsubsection{Optical Strain}

\textbf{Research gap.} The evidence for strain comes from hand-crafted descriptors evaluated under a protocol that is not subject-independent~\cite{liong2016}, and the strain magnitude of STSTNet cannot be determined unambiguously from its published equation~\cite{liong2019b}.

\textbf{Limitation.} The implemented magnitude counts shear once rather than twice, weighting it differently from the Frobenius norm used by Liong et al.~\cite{liong2014a,liong2016}, and no dedicated strain filter is applied~\cite{liong2014a}. Strain is present in all twelve configurations, so its own contribution is not measured.

\subsubsection{Convolutional Spatial Stem}

\textbf{Research gap.} The reviewed studies add modules to a backbone that remains in place~\cite{xia2020b}; none removes the learned spatial stem and passes the motion representation directly to the classifier, so its contribution is unknown. Deeper models and higher input resolutions reduce performance on composite micro-expression data~\cite{xia2020b}.

\textbf{Limitation.} The stem contains only spatial kernels, so the ablation tests a learned spatial stem rather than spatio-temporal convolution~\cite{ji2013,tran2015}. Its input resolution of $224\times224$ exceeds the $100\times100$ threshold reported by Xia et al.~\cite{xia2020b}, so any effect measured for the stem may reflect resolution as much as architecture.

\subsubsection{Parameter-Free Attention}

\textbf{Research gap.} SimAM has been evaluated on natural images with learned appearance features~\cite{yang2021}; its suitability for flow-derived features, and whether its statistics should be computed per frame or per clip, are untested. Its regularisation constant was selected separately for each image dataset~\cite{yang2021}.

\textbf{Limitation.} SimAM statistics are computed over the whole clip rather than per frame, and $\lambda = 10^{-4}$ is not re-tuned. Because SimAM requires the convolutional stem, its effect is measured over only four matched pairs and is conditional on the presence of the stem.

\subsection{Temporal Modelling}
\label{sec:gaps-temporal}

\subsubsection{Temporal Normalisation}

\textbf{Research gap.} Sequence length is usually fixed at ten frames by convention~\cite{liX2013,lu2015}, and the studies that vary it report different optima~\cite{liX2018,ben2021,xu2017}. No reviewed study compares interpolation methods for a motion-based representation, in which synthesised frames would produce synthesised displacements; Xu et al. prefer linear interpolation for motion features but do not compare it with TIM~\cite{xu2017}.

\textbf{Limitation.} Temporal normalisation is implemented as uniform sampling of existing frames rather than the Temporal Interpolation Model~\cite{pfister2011}, so short clips cannot be lengthened and three clips repeat frames. The length of 33 frames is not varied. Sampling reduces the effective frame rate to approximately 100~fps at the median clip length and 52~fps for the longest clip (Section~\ref{sec:three-channel-tensor}).

\subsubsection{Temporal Transformer}

\textbf{Research gap.} The strongest transformer-based results combine a pre-trained spatial transformer with an LSTM aggregator~\cite{zhang2022}, so the benefit of a temporal transformer trained from scratch on one small dataset has not been isolated. Transformers lack the inductive biases of convolution and can perform poorly when trained on insufficient data~\cite{dosovitskiy2021}.

\textbf{Limitation.} The implemented encoder differs from SLSTT in six respects (\autoref{tab:slstt-divergences}). The ablation compares the complete encoder, including positional encoding, feed-forward layers and normalisation, with mean pooling; it does not isolate self-attention or compare it with another learned temporal model.

\subsection{Training and Evaluation}
\label{sec:gaps-training}

\textbf{Research gap.} How imbalance corrections interact is examined by only one reviewed study~\cite{xia2020a}, the focusing parameter of focal loss has not been validated for this task~\cite{zhaoS2021}, and label smoothing~\cite{szegedy2016} has not been examined for micro-expression recognition.

\textbf{Limitation.} The imbalance treatment, $\gamma$ and label smoothing are identical in all configurations, so their effects are not measured. All results come from one dataset, one random seed and a checkpoint selection without an inner validation split (Sections~\ref{sec:design-limitation} and~\ref{sec:checkpoint-selection}), and no result extends to other datasets, populations or recording conditions (Section~\ref{sec:sibling-corpora}).


% =============================================================================
\section{Summary}
\label{sec:ch4-conclusion}
% =============================================================================

This chapter described the factorial ablation of four pipeline components on CASME~II. It presented the 156-clip working set and its 25 leave-one-subject-out folds, the twelve valid configurations and the twenty matched pairs used to compare them, and the shared model architecture, training procedure and evaluation protocol, with pooled macro F1 as the primary metric. It closed with the research gaps and limitations of each motion-modelling and temporal-modelling component.

Chapter~\ref{ch:results} reports the results of the ablation: the performance of the twelve configurations and the contribution of each component, measured over its matched pairs.
. This file
% is the single source for Chapter 4; chapter4.tex and the chapter4_short_half
% files are superseded and share labels with it, so never include them.
%
% Rewritten 23 September 2026:
%   * Subsections consolidated; merged subsections keep every earlier label
%     (stacked \label commands), so references from other chapters resolve.
%   * Removed: "What is stored" and "Earlier evaluations" (labels
%     sec:what-is-stored, sec:excluded-evaluations, tab:five-evaluations),
%     all VRAM figures, the Computational Cost section (removed at the
%     author's request), and every code file, script and class name.
%   * The pipeline overview (formerly its own section) is now Section 4.4.1.
%   * New Section 4.8 combines the research gap and the limitations of each
%     component, grouped into motion modelling and temporal modelling.
%   * Figures fig4_1 ... fig4_4 are generated by tools/thesis_ch4_figures.py
%     from the values stated in this chapter.
%
% Requires in the preamble: amsmath, amssymb (\checkmark), graphicx, tikz with
%   \usetikzlibrary{arrows.meta, positioning, fit, backgrounds}
% =============================================================================

\chapter{Methodology}
\label{ch:methodology}

This chapter describes the study as it was carried out. Section~\ref{sec:research-approach} presents the experimental design, Section~\ref{sec:corpus-preparation} the dataset and its preparation, and Section~\ref{sec:ablation-matrix} the twelve configurations and the matched pairs used to compare them. Section~\ref{sec:architecture} describes the model architecture, Section~\ref{sec:training-procedure} the training procedure, and Section~\ref{sec:evaluation-protocol} the evaluation protocol. Section~\ref{sec:gaps-limitations} summarises the research gaps and limitations of each component, grouped into motion modelling and temporal modelling.


% =============================================================================
\section{Research Approach and Experimental Design}
\label{sec:research-approach}
% =============================================================================

\subsection{Factorial Ablation Design}
\label{sec:study-form}

This thesis does not propose and evaluate a single architecture. The study is designed as a factorial ablation: four components of a micro-expression recognition pipeline are each switched on or off, every valid combination is trained and evaluated under an identical protocol, and the contribution of each component is estimated from pairs of configurations that differ in that component alone. The object of study is therefore the set of configurations, and the question is not how well the pipeline performs but which of its stages are responsible for that performance, and by how much.

This design addresses the research gap of Section~\ref{sec:research-gap}. In the reviewed literature, components are varied one at a time within a single proposed system, for example the magnification factor~\cite{liX2018}, individual modules of a recurrent convolutional network~\cite{xia2020b} or the temporal aggregator of a transformer~\cite{zhang2022}. No reviewed study varies several components jointly under one protocol with all other factors held identical.

\subsection{Ablated Components}
\label{sec:four-variables}

The design space is defined by four components, each of which is either fully present or absent:

\begin{itemize}
	\item \textbf{Eulerian video magnification (EVM)}~\cite{wu2012}, applied during preprocessing before the motion representation is computed (Section~\ref{sec:motion-magnification}).
	\item \textbf{A convolutional spatial stem}, with one stream per motion channel (Section~\ref{sec:conv3d}).
	\item \textbf{SimAM}~\cite{yang2021}, a parameter-free attention module within the convolutional streams (Section~\ref{sec:simam}).
	\item \textbf{A temporal transformer encoder}~\cite{vaswani2017}, applied to the sequence of frame features (Section~\ref{sec:self-attention}).
\end{itemize}

The first three act on the motion representation and are grouped as \textit{motion modelling}; the transformer models the evolution of that representation over time and constitutes \textit{temporal modelling}. A disabled component is replaced by a minimal operation (Section~\ref{sec:fallbacks-config}) rather than by an alternative component of comparable capacity.

\subsection{Design Limitations}
\label{sec:design-limitation}
\label{sec:what-is-stored}

Every configuration is trained once, under a single fixed random seed (42). No configuration is retrained with other seeds, so no variance across seeds is estimated. Every difference between two configurations, including the matched-pair effects of Section~\ref{sec:matched-pairs}, is therefore a difference between two point estimates without a confidence interval, and it could change under a different seed. In addition, per-clip predictions were not retained; only aggregate metrics and confusion matrices were stored. A paired significance test such as McNemar's test~\cite{mcnemar1947}, which requires the prediction for each clip, therefore cannot be computed for any pair of configurations.


% =============================================================================
\section{Dataset and Preparation}
\label{sec:corpus-preparation}
% =============================================================================

\subsection{Data Selection}
\label{sec:corpus-used}
\label{sec:255-to-156}

All experiments use CASME~II~\cite{yan2014} (Section~\ref{sec:corpus}); no other dataset contributes to training or evaluation, so every result is a within-dataset result. The released coding contains 255 clips coded as micro-expressions with extractable frames (Section~\ref{sec:clip-selection}). The \texttt{Others} category (99 clips) is excluded before training. It is a residual category for clips that do not fit the other labels, and a class defined by exclusion rather than content cannot provide a meaningful training signal. Figure~\ref{fig:class-distribution} shows the composition of the 255 clips and the resulting working set of 156 clips.

\begin{figure}[htbp]
	\centering
	\includegraphics[width=0.85\textwidth]{fig4_1_class_distribution.png}
	\caption[Composition of the CASME II working set]{Composition of the 255 CASME~II micro-expression clips by emotion label, grouped into the three classes used in this thesis. The \texttt{Others} category is excluded, leaving a working set of 156 clips. Counts are taken from the released coding (Section~\ref{sec:label-taxonomy}).}
	\label{fig:class-distribution}
\end{figure}

\subsection{Class Grouping}
\label{sec:three-class-grouping}

Disgust, repression, sadness and fear are grouped into Negative, happiness forms Positive and surprise forms Surprise, following the valence-based convention of SMIC and MEGC 2019~\cite{liX2013,see2019} (Section~\ref{sec:label-taxonomy}). The resulting class distribution is 99 : 32 : 25, approximately 4 : 1.3 : 1. All results are reported for these three classes.

\subsection{Subjects and Folds}
\label{sec:subjects-loso}

The 255 clips come from 26 subjects. Subject 18 contributes only \texttt{Others} clips and is removed together with that category, so the working set contains 25 subjects, and leave-one-subject-out cross-validation produces \textbf{25 folds} rather than 26. Because subjects contribute unequal numbers of clips, the folds differ considerably in size and class composition (\autoref{fig:fold-composition}).

\begin{figure}[tbp]
	\centering
	\includegraphics[width=\textwidth]{fig3_2_fold_composition.png}
	\caption[Class composition of the 25 leave-one-subject-out folds]{Class composition of the 25 leave-one-subject-out folds. Top: clips held out in each fold, by class, with the number of classes present above each bar. Bottom: the number of folds containing one, two or three classes, and the per-fold macro-F1 ceiling this implies.}
	\label{fig:fold-composition}
\end{figure}

\subsection{Input Preparation}
\label{sec:network-preparation}

Each clip is converted into a motion tensor of shape $(3, 32, 224, 224)$, whose three channels hold horizontal optical flow, vertical optical flow and optical strain magnitude over 32 frame pairs (Sections~\ref{sec:recording-to-input} and~\ref{sec:video-to-motion}). Table~\ref{tab:preprocessing-settings} lists the settings used. None of them is varied by the ablation, so every configuration receives inputs of the same shape.

\begin{table}[htbp]
	\centering
	\small
	\begin{tabular}{@{}l l@{}}
		\hline
		\textbf{Setting} & \textbf{Value} \\
		\hline
		Frames sampled per clip        & 33, uniformly between onset and offset (32 flow pairs) \\
		Frame size                     & $224\times224$ pixels, greyscale \\
		Magnification factor $\alpha$  & 10 \\
		Magnification band             & 5--25 Hz at a nominal 200 fps \\
		Laplacian pyramid levels       & 4 \\
		Magnified output               & clipped to $[0,255]$ and converted to 8-bit \\
		Optical flow estimator         & Farneb\"ack~\cite{farneback2003} \\
		Channel normalisation          & min--max to $[0,1]$ per clip \\
		\hline
	\end{tabular}
	\caption[Preprocessing settings]{Preprocessing settings, fixed across all configurations. Section~\ref{sec:evm-departures} lists how the magnification settings depart from the reference formulation.}
	\label{tab:preprocessing-settings}
\end{table}


% =============================================================================
\section{Ablation Configurations}
\label{sec:ablation-matrix}
% =============================================================================

\subsection{Valid Configurations}
\label{sec:twelve-configs}
\label{sec:twelve-trained}
\label{sec:excluded-cells}
\label{sec:naming}

The four binary components define $2^4 = 16$ combinations (\autoref{fig:ablation-grid}). Four of them enable SimAM without the convolutional stem. Because SimAM re-weights the activations of a convolutional feature map, it has nothing to act on without the stem; these combinations are therefore not weaker models but ill-formed ones, and they are not trained. The remaining twelve configurations are listed in \autoref{tab:twelve-configs}. Configurations 1 to 8 were defined first and configurations 9, 12, 13 and 16 were added to complete the design, so the numbering omits the four excluded combinations (10, 11, 14 and 15). Configuration $n$ is denoted C$n$ in tables and figures.

\begin{figure}[htbp]
	\centering
	\includegraphics[width=0.8\textwidth]{fig4_2_ablation_grid.png}
	\caption[The factorial ablation design]{The $2^4$ factorial design. Each cell is one combination of the four components; twelve are trained and the four that enable SimAM without the convolutional stem are excluded.}
	\label{fig:ablation-grid}
\end{figure}

\begin{table}[htbp]
	\centering
	\small
	\begin{tabular}{@{}l c c c c@{}}
		\hline
		\textbf{Configuration} & \textbf{EVM} & \textbf{SimAM} & \textbf{CNN} & \textbf{TF} \\
		\hline
		C1   & --         & --         & --         & -- \\
		C2   & --         & --         & --         & \checkmark \\
		C3   & --         & --         & \checkmark & -- \\
		C4   & \checkmark & --         & --         & -- \\
		C5   & --         & \checkmark & \checkmark & -- \\
		C6   & --         & \checkmark & \checkmark & \checkmark \\
		C7   & \checkmark & --         & \checkmark & \checkmark \\
		C8   & \checkmark & \checkmark & \checkmark & \checkmark \\
		C9   & --         & --         & \checkmark & \checkmark \\
		C12  & \checkmark & --         & --         & \checkmark \\
		C13  & \checkmark & --         & \checkmark & -- \\
		C16  & \checkmark & \checkmark & \checkmark & -- \\
		\hline
	\end{tabular}
	\caption[The twelve valid ablation configurations]{The twelve valid ablation configurations. EVM: Eulerian video magnification; CNN: convolutional stem; TF: temporal transformer.}
	\label{tab:twelve-configs}
\end{table}

\subsection{Baseline and Control}
\label{sec:baseline-control}

Two configurations serve different purposes and should not be confused. C4 is the baseline of the study: it applies magnification and no learned component, and the contributions of the learned components are ultimately read against it. C1 is not a second baseline but the control for magnification: it is identical to C4 except that magnification is disabled. C8 enables all four components but is otherwise one of twelve configurations and is not privileged in the analysis.

\subsection{Matched Pairs}
\label{sec:matched-pairs}
\label{sec:matched-pairs-table}

A matched pair consists of two configurations that differ in exactly one component. The difference between their scores estimates the contribution of that component with all others held fixed, and the effect reported for a component is the mean over all of its pairs. Magnification and the transformer can be toggled in each of the six valid combinations of the other three components, giving six pairs each. A SimAM pair requires the stem to be present, and a stem pair requires SimAM to be absent, so these two components have four pairs each. Table~\ref{tab:matched-pairs} lists all twenty pairs. Because an effect averaged over four pairs is a noisier estimate than one averaged over six, effects with different pair counts should be compared with care.

\begin{table}[htbp]
	\centering
	\small
	\begin{tabular}{@{}l l r@{}}
		\hline
		\textbf{Component} & \textbf{Pairs (off $\rightarrow$ on)} & \textbf{Count} \\
		\hline
		Transformer & 1$\rightarrow$2, 4$\rightarrow$12, 3$\rightarrow$9, 13$\rightarrow$7, 5$\rightarrow$6, 16$\rightarrow$8 & 6 \\
		EVM         & 1$\rightarrow$4, 3$\rightarrow$13, 9$\rightarrow$7, 6$\rightarrow$8, 5$\rightarrow$16, 2$\rightarrow$12 & 6 \\
		SimAM       & 3$\rightarrow$5, 9$\rightarrow$6, 7$\rightarrow$8, 13$\rightarrow$16 & 4 \\
		CNN stem    & 1$\rightarrow$3, 2$\rightarrow$9, 4$\rightarrow$13, 12$\rightarrow$7 & 4 \\
		\hline
	\end{tabular}
	\caption[Matched pairs by component]{Matched pairs by component, identified by configuration number.}
	\label{tab:matched-pairs}
\end{table}


% =============================================================================
\section{Model Architecture}
\label{sec:architecture}
% =============================================================================

All twelve configurations are branches of a single model in which each component is executed or skipped according to switches fixed at construction; magnification is applied during preprocessing and is not part of the network. Because every configuration shares the same implementation, initialisation and layer order, two configurations that differ in one switch differ in nothing else.

\subsection{Pipeline Overview}
\label{sec:pipeline-overview}
\label{sec:pipeline-in-words}
\label{sec:pipeline-shapes-discussion}

The pipeline consists of two stages that are executed a different number of times (\autoref{fig:pipeline-overview}). \textbf{Stage~1 (preprocessing)} converts each clip into one motion tensor and runs once per clip. \textbf{Stage~2 (training and evaluation)} loads these tensors and fits a model; it runs once for every configuration and every held-out subject. Because Stage~1 is completed before any model is trained, all configurations are compared on identical inputs.

% -----------------------------------------------------------------------------
\begin{figure}[htbp]
	\centering
	\begin{tikzpicture}[
		font=\scriptsize,
		node distance=3.2mm,
		box/.style={
			draw, rounded corners=2pt, align=center,
			text width=34mm, inner sep=3pt,
			fill=black!4},
		switch/.style={box, dashed, fill=black!10},
		term/.style={box, fill=black!14, font=\scriptsize\bfseries},
		arr/.style={-{Stealth[length=1.6mm]}, semithick}
	]

	% ---------------- Stage 1: preprocessing ----------------
	\node[term] (a1) {CASME II dataset\\\normalfont 200\,fps, FACS-coded};
	\node[box, below=of a1] (a2)
		{Clip selection\\\normalfont 156 clips, 25 subjects};
	\node[box, below=of a2] (a3)
		{Frame preparation\\\normalfont greyscale, $224\times224$};
	\node[box, below=of a3] (a4)
		{Uniform sampling\\\normalfont 33 frames per clip};
	\node[switch, below=of a4] (a5)
		{Motion magnification\\\normalfont Eulerian, $\alpha=10$};
	\node[box, below=of a5] (a6)
		{Optical flow\\\normalfont Farneb\"ack, 32 pairs};
	\node[box, below=of a6] (a7)
		{Optical strain\\\normalfont spatial gradients};
	\node[box, below=of a7] (a8)
		{Stack and normalise\\\normalfont $[u,v,\varepsilon]$ to $[0,1]$};
	\node[term, below=of a8] (a9)
		{Saved motion tensor\\\normalfont $(3,32,224,224)$};

	\foreach \i/\j in {a1/a2, a2/a3, a3/a4, a4/a5, a5/a6, a6/a7, a7/a8, a8/a9}
		{\draw[arr] (\i) -- (\j);}

	% ---------------- Stage 2: training and evaluation ----------------
	\node[box, right=20mm of a2] (b1)
		{Load and standardise\\\normalfont flip augmentation};
	\node[switch, below=of b1] (b2)
		{Convolutional stem\\\normalfont three streams};
	\node[switch, below=of b2] (b3)
		{SimAM\\\normalfont parameter-free gate};
	\node[box, below=of b3] (b4)
		{Pool to a sequence\\\normalfont 32 steps, 96 features};
	\node[switch, below=of b4] (b5)
		{Temporal transformer\\\normalfont 4 layers, 8 heads};
	\node[box, below=of b5] (b6)
		{Classifier head\\\normalfont three logits};
	\node[box, below=of b6] (b7)
		{Leave-one-subject-out\\\normalfont 25 folds, counts pooled};
	\node[term, below=of b7] (b8)
		{Pooled macro F1\\\normalfont the primary metric};

	\foreach \i/\j in {b1/b2, b2/b3, b3/b4, b4/b5, b5/b6, b6/b7, b7/b8}
		{\draw[arr] (\i) -- (\j);}

	% ---------------- the hand-off ----------------
	\draw[arr] (a9.east) -- ++(8mm,0) |- (b1.west);

	% ---------------- stage brackets ----------------
	\begin{scope}[on background layer]
		\node[draw=black!45, dashed, rounded corners=3pt, inner sep=4.5mm,
		      fit=(a1)(a9)] (s1) {};
		\node[draw=black!45, dashed, rounded corners=3pt, inner sep=4.5mm,
		      fit=(b1)(b8)] (s2) {};
	\end{scope}
	\node[above=0.5mm of s1, font=\scriptsize\bfseries] {Stage 1: preprocessing};
	\node[above=0.5mm of s2, font=\scriptsize\bfseries] {Stage 2: training and evaluation};

	\end{tikzpicture}
	\caption[The data pipeline, end to end]{The pipeline from dataset to score. Stage~1 runs once per clip; Stage~2 runs once per configuration and fold. Solid boxes are always executed. \textbf{Dashed boxes are the four components switched on and off in the ablation}; each is replaced by a defined operation when disabled.}
	\label{fig:pipeline-overview}
\end{figure}
% -----------------------------------------------------------------------------

In Stage~1, each selected clip is loaded in greyscale, resized to $224\times224$ pixels and sampled at 33 evenly spaced frames between onset and offset. Magnification is applied at this point if enabled. Optical flow is then computed for the 32 adjacent frame pairs, optical strain is derived from its spatial gradients, and the three fields are stacked and normalised into one tensor. The network therefore never receives pixel intensities, only displacement and deformation fields. In Stage~2, each tensor is standardised per channel, optionally flipped for augmentation, and passed through the convolutional stem, the SimAM gate, a pooling step that removes the spatial dimensions, the temporal encoder and the classification head. Table~\ref{tab:pipeline-shapes} lists the representation of the data after each step.

\begin{table}[htbp]
	\centering
	\small
	\begin{tabular}{@{}lll@{}}
		\hline
		\textbf{After this step} & \textbf{The data is} & \textbf{Shape} \\
		\hline
		Clip selection        & a set of clips                & 156 clips, 25 subjects \\
		Frame preparation     & greyscale frames              & $N \times 224 \times 224$ \\
		Uniform sampling      & a fixed-length sequence       & $33 \times 224 \times 224$ \\
		Magnification         & the same, amplified           & $33 \times 224 \times 224$ \\
		Flow and strain       & three motion fields per pair  & $3 \times 32 \times 224 \times 224$ \\
		Pooling to a sequence & a sequence of feature vectors & $32 \times 96$ \\
		Temporal encoding     & one vector per clip           & $96$ \\
		Classifier head       & class scores                  & $3$ \\
		Fold aggregation      & pooled confusion counts       & $3 \times 3$ \\
		\hline
	\end{tabular}
	\caption[What the data is at each pipeline step]{The data representation after each pipeline step, where $N$ is the number of frames between onset and offset. The temporal axis falls from 33 frames to 32 pairs because flow is computed between \emph{adjacent} frames, and the spatial axes are removed at the pooling step rather than at the classifier.}
	\label{tab:pipeline-shapes}
\end{table}

\subsection{Convolutional Stem}
\label{sec:input-streams}
\label{sec:conv-stem}

Each of the three motion channels is processed by its own convolutional stream with unshared weights. Table~\ref{tab:conv-stem} lists the layers of one stream, which use batch normalisation~\cite{ioffe2015}, ReLU activations~\cite{nair2010} and channel-wise dropout~\cite{tompson2015}. No kernel spans the temporal axis (Section~\ref{sec:conv3d}), so the 32 time steps pass through unchanged while height and width are halved once. The outputs of the three streams are concatenated along the channel dimension, giving a feature map of $96\times32\times112\times112$ per clip.

\begin{table}[htbp]
	\centering
	\small
	\begin{tabular}{@{}l l l@{}}
		\hline
		\textbf{Layer} & \textbf{Configuration} & \textbf{Output (per stream)} \\
		\hline
		3D convolution          & 16 filters, kernel $(1,3,3)$, padding $(0,1,1)$, no bias & $16\times32\times224\times224$ \\
		Normalisation and activation & batch normalisation, ReLU, 3D dropout ($p=0.3$) & $16\times32\times224\times224$ \\
		3D convolution          & 32 filters, kernel $(1,3,3)$, padding $(0,1,1)$, no bias & $32\times32\times224\times224$ \\
		Normalisation and activation & batch normalisation, ReLU, 3D dropout ($p=0.3$) & $32\times32\times224\times224$ \\
		3D max pooling          & kernel and stride $(1,2,2)$ & $32\times32\times112\times112$ \\
		\hline
	\end{tabular}
	\caption[Layers of one convolutional stream]{Layers of one convolutional stream. Shapes are given as channels $\times$ time $\times$ height $\times$ width.}
	\label{tab:conv-stem}
\end{table}

\subsection{SimAM}
\label{sec:simam-placement}

When enabled, SimAM~\cite{yang2021} is applied separately to the output of each convolutional stream, before the streams are concatenated. It follows \eqref{eq:simam-stats} and \eqref{eq:simam-gate}, with the mean and variance computed over the time, height and width of each channel (divisor $THW-1$) and $\lambda = 10^{-4}$. It adds no learnable parameters. SimAM is enabled only together with the convolutional stem.

\subsection{Temporal Transformer}
\label{sec:transformer-config}

The concatenated feature map is reduced by adaptive average pooling over the two spatial dimensions to a sequence of 32 time steps, each a 96-dimensional feature vector. Fixed sinusoidal positional encodings~\cite{vaswani2017} are added to this sequence:
\begin{equation}
	\mathrm{PE}(p, 2i) = \sin\!\left(\frac{p}{10000^{2i/96}}\right), \qquad \mathrm{PE}(p, 2i+1) = \cos\!\left(\frac{p}{10000^{2i/96}}\right),
	\label{eq:positional-encoding}
\end{equation}
where $\mathrm{PE}(p, j)$ is the value added to feature $j$ of the token at time step $p \in \{0, \dots, 31\}$, $i$ indexes the pairs of feature dimensions, and 96 is the feature dimension. The encoding is fixed rather than learned, and dropout with probability 0.1 is applied after it is added. The sequence is then processed by a transformer encoder with four layers, eight attention heads, model dimension 96, feed-forward dimension 256, dropout 0.1, GELU activation~\cite{hendrycks2016} and pre-norm layer ordering~\cite{xiong2020}, followed by a final layer normalisation~\cite{ba2016}. The output sequence is averaged over time, giving one 96-dimensional vector per clip.

\subsection{Classification Head and Replacement Operations}
\label{sec:classifier-head}
\label{sec:fallbacks-config}

Every configuration ends in the same classification head: layer normalisation, dropout with probability 0.3, and a linear layer mapping the 96 features to three class logits. A disabled component is replaced by a minimal operation rather than removed, so that the configuration remains a fair control (\autoref{tab:ch2-component-controls}). Without the convolutional stem, each frame of the input tensor is average-pooled to a $4\times4$ grid per channel, and the resulting 48 values are projected to 96 features by a linear layer with 4,704 parameters. Without the transformer, the sequence is averaged over time; this operation has no parameters and no notion of temporal order.

\subsection{Parameter Counts}
\label{sec:parameter-counts}
\label{sec:weight-init}

The convolutional stem has 14,544 parameters across its three streams, the transformer encoder 348,736, the replacement projection 4,704 and the classification head 483; SimAM and magnification add none. Figure~\ref{fig:parameter-counts} shows the resulting totals for the four architectural combinations. In a configuration with all components, the transformer holds approximately 96\,\% of the parameters. All configurations use the same initialisation: Kaiming-normal initialisation for convolutional weights, Xavier-uniform initialisation for linear weights, unit scale and zero shift for normalisation layers, and zero biases.

\begin{figure}[htbp]
	\centering
	\includegraphics[width=0.85\textwidth]{fig4_3_parameter_counts.png}
	\caption[Parameter totals by architectural combination]{Trainable parameters of the four architectural combinations (logarithmic scale), with the contribution of each component.}
	\label{fig:parameter-counts}
\end{figure}


% =============================================================================
\section{Training Procedure}
\label{sec:training-procedure}
% =============================================================================

All twelve configurations are trained with the same procedure, summarised in \autoref{tab:training-settings}; the underlying methods are described in Section~\ref{sec:scarcity-skew}.

\begin{table}[htbp]
	\centering
	\small
	\begin{tabular}{@{}l p{0.62\linewidth}@{}}
		\hline
		\textbf{Setting} & \textbf{Value} \\
		\hline
		Loss            & focal loss~\cite{lin2017} with $\gamma = 2$~\cite{zhaoS2021} and label smoothing~\cite{szegedy2016} of 0.05; no class weights \\
		Sampling        & class-balanced, weight $1/n_c$ per example, with replacement \\
		Augmentation    & random horizontal flip ($p = 0.5$) with negation of the horizontal flow \\
		Optimiser       & AdamW~\cite{loshchilov2019}, learning rate $10^{-4}$, weight decay $10^{-4}$ \\
		Learning rate   & linear warmup over 5 epochs from $10^{-5}$, then cosine annealing~\cite{loshchilov2017} to $10^{-7}$ \\
		Epochs          & 50, without early stopping \\
		Batch size      & 8 \\
		Gradient clipping & global norm 1.0 \\
		Precision       & mixed precision with gradient scaling~\cite{micikevicius2018} \\
		Random seed     & 42, reset before every configuration and fold \\
		\hline
	\end{tabular}
	\caption[Training settings]{Training settings, identical for all configurations and folds.}
	\label{tab:training-settings}
\end{table}

\subsection{Loss and Class Weighting}
\label{sec:loss-config}
\label{sec:class-weighting-rule}
\label{sec:sampler-config}

The loss can apply per-class weights ($\alpha_t$ in \eqref{eq:focal-loss}), but they are enabled only when the class-balanced sampler is inactive. Because the sampler is active in every reported run, the loss carries \textbf{no class weights} (Section~\ref{sec:combining-both}), and the frequency correction comes entirely from the sampler, which gives each training example a weight inversely proportional to the size of its class in the training split of the current fold \eqref{eq:sampler-weight}. The focal term $(1-p_t)^\gamma$ and label smoothing remain active; neither depends on class frequency.

The rule was introduced after two failures in opposite directions, both observed during the development of this study. In an earlier version, the loss carried inverse-frequency class weights and no sampler was used; under this single correction, the proposed configuration stopped predicting the \textbf{Positive} class altogether. A later version added the sampler while keeping the loss weights, so both corrections operated together; the failure then reversed, and configurations stopped predicting the \textbf{Negative} majority class, one of them predicting a single class for every clip. With the rule in place, no configuration in the reported study abandons any class. The only reviewed study that combines two corrections reports that the balanced loss adds only a slight improvement, not that the combination reverses its effect~\cite{xia2020a}. The two failures, however, occurred in implementation versions that differ from the reported run in more than the correction rule, so \textbf{the causal effect of the rule is not isolated}.

\subsection{Optimisation}
\label{sec:optimiser-config}
\label{sec:batch-size}
\label{sec:numerical-safeguards}

Figure~\ref{fig:lr-schedule} shows the learning-rate schedule. The learning rate rises linearly from one tenth of its target value over the first five epochs and then follows a cosine decay to $10^{-7}$; the schedule is updated once per epoch. Every configuration and fold is trained for the full 50 epochs. Gradients are clipped to a global norm of 1.0 after being unscaled from mixed precision, and an update is skipped if any gradient contains a non-finite value.

\begin{figure}[htbp]
	\centering
	\includegraphics[width=0.8\textwidth]{fig4_4_lr_schedule.png}
	\caption[Learning-rate schedule]{Learning rate per epoch: linear warmup over five epochs followed by cosine annealing.}
	\label{fig:lr-schedule}
\end{figure}

\subsection{Reproducibility}
\label{sec:seeding}
\label{sec:no-hyperparameter-search}

The random seed is reset before every configuration and every fold, so the sampling order, weight initialisation and all other stochastic steps start from the same state each time a fold is trained. No hyperparameter search was performed: all training settings are fixed defaults, identical across the twelve configurations and the 25 folds. Every comparison is therefore between architectures trained with one common, untuned recipe, not between architectures each tuned to its own advantage.


% =============================================================================
\section{Evaluation Protocol}
\label{sec:evaluation-protocol}
% =============================================================================

All configurations are evaluated with the same fold structure, checkpoint selection rule and aggregation across folds. The underlying metrics are defined in Section~\ref{sec:evaluation}.

\subsection{Leave-One-Subject-Out Folds}
\label{sec:loso-folds}

Evaluation uses leave-one-subject-out cross-validation~\cite{yan2014,see2019} over the 25 subjects of the working set, giving \textbf{25 folds}. Each fold holds out all clips of one subject and trains on the clips of the remaining 24 subjects. All 25 folds are run for every configuration.

\subsection{Checkpoint Selection}
\label{sec:checkpoint-selection}

Within each fold, the macro F1 on the held-out subject is computed after every epoch, and the weights of the best epoch are restored and scored on the same held-out subject. \textbf{There is no inner validation split}: the held-out clips are used both to select the epoch and to score it. A proper nested procedure would select the epoch on a validation subset of the training subjects and reserve the held-out subject for the final evaluation. The reported scores are therefore optimistically biased by an unknown amount~\cite{varma2006}. The bias is likely to be largest for the smallest folds, where the choice among 50 epochs is made on only a few clips.

\subsection{Aggregation and Primary Metric}
\label{sec:aggregation}
\label{sec:primary-metric}
\label{sec:mean-of-folds-rejected}

The confusion matrices of the 25 folds are summed into one pooled $3\times3$ matrix, from which the per-class precision, recall and F1 are computed. The primary metric is \textbf{pooled macro F1}, the mean of the three pooled per-class F1 scores, which is identical in construction to the Unweighted F1 of MEGC 2019~\cite{see2019}. Pooled accuracy is reported alongside it.

The alternative, averaging the macro F1 of each fold, is rejected because of the fold composition. Of the 25 folds, 10 contain clips of only one class, 8 contain two and 7 contain all three. A class absent from a fold contributes an F1 of zero to that fold, so a single-class fold can score at most $1/3$ and a two-class fold at most $2/3$. The mean-of-folds macro F1 is therefore bounded by
\begin{equation}
	\overline{F1}_{\max} = \frac{10 \cdot \tfrac13 + 8 \cdot \tfrac23 + 7 \cdot 1}{25} = \frac{3.333\ldots + 5.333\ldots + 7}{25} \approx 0.627,
	\label{eq:fold-ceiling}
\end{equation}
where $\overline{F1}_{\max}$ is the highest mean-of-folds macro F1 that any configuration can reach, the numerators weight each fold type by its maximum score, and 25 is the number of folds. No configuration can exceed this ceiling, however accurate it is. Pooling avoids the ceiling, because true positives, false positives and false negatives are accumulated over all 156 clips before any F1 is computed. Fold composition also inflates mean-of-folds accuracy: for the configuration with all four components, it exceeds pooled accuracy by 6.3 percentage points.


% =============================================================================
\section{Research Gaps and Limitations}
\label{sec:gaps-limitations}
% =============================================================================

This section combines, for each component, the research gap left by the reviewed literature (Chapter~\ref{ch:literature}) and the limitations of its implementation in this thesis. The components are grouped into motion modelling, which covers the representation and spatial processing of facial motion, and temporal modelling, which covers the sampling of the sequence and the modelling of its evolution over time. Limitations of the training and evaluation procedure as a whole are listed last.

\subsection{Motion Modelling}
\label{sec:gaps-motion}

\subsubsection{Eulerian Video Magnification}

\textbf{Research gap.} Magnification improves recognition in the reviewed studies, but every evaluation places it in front of an appearance-based representation such as LBP, HOG or pre-trained face features~\cite{liX2018,bai2021,liY2018,liY2021}. No study measures its contribution when the classifier receives an explicit motion field. The magnification factor has an interior optimum that can be found only by search~\cite{liX2018}, and magnification also amplifies unwanted motion such as head movement~\cite{liX2018}.

\textbf{Limitation.} The factor is fixed at $\alpha = 10$, inside the reported optimum, but is not varied. The pass-band, the filter and the order of magnification and sampling depart from the reference formulation (\autoref{tab:ch2-evm-departures}); in particular, the realised pass-band depends on the clip length and is approximately 2.5--12.5~Hz at the median length. No face alignment precedes magnification, so head motion within the band is amplified together with facial movement.

\subsubsection{Optical Flow}

\textbf{Research gap.} Optical flow is the preferred input of the strongest benchmark systems~\cite{see2019}, but the choice of estimator is not settled: Farneb\"ack and TV-L1 are both used~\cite{zhaoS2021,liong2019a}, and no controlled comparison of their recognition performance exists. Whether a single onset-to-apex flow field~\cite{liong2019a,liong2019b} or a sequence of flow fields~\cite{zhaoS2021} is preferable has likewise not been measured.

\textbf{Limitation.} Only the Farneb\"ack estimator~\cite{farneback2003} is used, and flow is computed between consecutive sampled frames rather than relative to the onset. No flow-based fine alignment is performed, although a small misalignment can considerably degrade micro-expression recognition~\cite{xu2017}.

\subsubsection{Optical Strain}

\textbf{Research gap.} The evidence for strain comes from hand-crafted descriptors evaluated under a protocol that is not subject-independent~\cite{liong2016}, and the strain magnitude of STSTNet cannot be determined unambiguously from its published equation~\cite{liong2019b}.

\textbf{Limitation.} The implemented magnitude counts shear once rather than twice, weighting it differently from the Frobenius norm used by Liong et al.~\cite{liong2014a,liong2016}, and no dedicated strain filter is applied~\cite{liong2014a}. Strain is present in all twelve configurations, so its own contribution is not measured.

\subsubsection{Convolutional Spatial Stem}

\textbf{Research gap.} The reviewed studies add modules to a backbone that remains in place~\cite{xia2020b}; none removes the learned spatial stem and passes the motion representation directly to the classifier, so its contribution is unknown. Deeper models and higher input resolutions reduce performance on composite micro-expression data~\cite{xia2020b}.

\textbf{Limitation.} The stem contains only spatial kernels, so the ablation tests a learned spatial stem rather than spatio-temporal convolution~\cite{ji2013,tran2015}. Its input resolution of $224\times224$ exceeds the $100\times100$ threshold reported by Xia et al.~\cite{xia2020b}, so any effect measured for the stem may reflect resolution as much as architecture.

\subsubsection{Parameter-Free Attention}

\textbf{Research gap.} SimAM has been evaluated on natural images with learned appearance features~\cite{yang2021}; its suitability for flow-derived features, and whether its statistics should be computed per frame or per clip, are untested. Its regularisation constant was selected separately for each image dataset~\cite{yang2021}.

\textbf{Limitation.} SimAM statistics are computed over the whole clip rather than per frame, and $\lambda = 10^{-4}$ is not re-tuned. Because SimAM requires the convolutional stem, its effect is measured over only four matched pairs and is conditional on the presence of the stem.

\subsection{Temporal Modelling}
\label{sec:gaps-temporal}

\subsubsection{Temporal Normalisation}

\textbf{Research gap.} Sequence length is usually fixed at ten frames by convention~\cite{liX2013,lu2015}, and the studies that vary it report different optima~\cite{liX2018,ben2021,xu2017}. No reviewed study compares interpolation methods for a motion-based representation, in which synthesised frames would produce synthesised displacements; Xu et al. prefer linear interpolation for motion features but do not compare it with TIM~\cite{xu2017}.

\textbf{Limitation.} Temporal normalisation is implemented as uniform sampling of existing frames rather than the Temporal Interpolation Model~\cite{pfister2011}, so short clips cannot be lengthened and three clips repeat frames. The length of 33 frames is not varied. Sampling reduces the effective frame rate to approximately 100~fps at the median clip length and 52~fps for the longest clip (Section~\ref{sec:three-channel-tensor}).

\subsubsection{Temporal Transformer}

\textbf{Research gap.} The strongest transformer-based results combine a pre-trained spatial transformer with an LSTM aggregator~\cite{zhang2022}, so the benefit of a temporal transformer trained from scratch on one small dataset has not been isolated. Transformers lack the inductive biases of convolution and can perform poorly when trained on insufficient data~\cite{dosovitskiy2021}.

\textbf{Limitation.} The implemented encoder differs from SLSTT in six respects (\autoref{tab:slstt-divergences}). The ablation compares the complete encoder, including positional encoding, feed-forward layers and normalisation, with mean pooling; it does not isolate self-attention or compare it with another learned temporal model.

\subsection{Training and Evaluation}
\label{sec:gaps-training}

\textbf{Research gap.} How imbalance corrections interact is examined by only one reviewed study~\cite{xia2020a}, the focusing parameter of focal loss has not been validated for this task~\cite{zhaoS2021}, and label smoothing~\cite{szegedy2016} has not been examined for micro-expression recognition.

\textbf{Limitation.} The imbalance treatment, $\gamma$ and label smoothing are identical in all configurations, so their effects are not measured. All results come from one dataset, one random seed and a checkpoint selection without an inner validation split (Sections~\ref{sec:design-limitation} and~\ref{sec:checkpoint-selection}), and no result extends to other datasets, populations or recording conditions (Section~\ref{sec:sibling-corpora}).


% =============================================================================
\section{Summary}
\label{sec:ch4-conclusion}
% =============================================================================

This chapter described the factorial ablation of four pipeline components on CASME~II. It presented the 156-clip working set and its 25 leave-one-subject-out folds, the twelve valid configurations and the twenty matched pairs used to compare them, and the shared model architecture, training procedure and evaluation protocol, with pooled macro F1 as the primary metric. It closed with the research gaps and limitations of each motion-modelling and temporal-modelling component.

Chapter~\ref{ch:results} reports the results of the ablation: the performance of the twelve configurations and the contribution of each component, measured over its matched pairs.
